% 제5장 논의 및 결론
% 1차 심사 반영: 구 4.7 종합 논의 + 4.8--4.10 시사점 + 4.11 한계 + 구 제5장 결론(요약/성과/기여/향후) 통합

\chapter{\vrev{논의 및 결론}}
\label{ch:discussion}

\vrev{%
본 연구에서는 제4장(연구 결과)에 대한 종합 논의를 제시하고,
\chg{\m{이론적·방법론적·실무적}} 시사점을 도출한다.
이어서 본 연구의 한계를 기술하고, \chg{연구 요약, 주요 성과, 학술적 기여}를 정리하며,
향후 연구 방향을 제시함으로써 본 연구를 마무리한다.
}

% ─── 5단계 5-1단계 (2026-04-22): Plan H 구조 재편 ───
% 9개 절 → 5개 절 통합, 본문 텍스트 무수정, wrapper로 깊이 강등
%
% 새 구조:
%   5.1 종합 논의 (현 5.1 그대로 — 5.1.2 이론적 기여, 5.1.3 연구의 한계 잔존, 5-2에서 정리)
%   5.2 학술적 기여 (wrapper: 현 5.2 + 5.3 + 5.8 흡수 + 5.2.3 실증적 기여 신설)
%   5.3 실무적 기여 (wrapper: 현 5.4 흡수, 5-2에서 3항으로 재분류)
%   5.4 한계점 (현 sec5.4_limitations.tex 그대로, 절 번호만 자동 5.4)
%   5.5 결론 및 향후 연구 (wrapper: 현 5.6 + 5.7 + 5.9 통합)

% 5-2단계: 신규 목차 구조 적용 (본문 텍스트 무수정, 파일 분할 + suppression wrapper)
% =====================================================================
% 5.1 종합 논의 — 5단계 5-2 (메인 wrapper, 5.1.2/5.1.3 분리됨)
% sec5.1_synthesis.tex의 분할본을 \input으로 결합. 본문 텍스트 무수정.
% =====================================================================



%% ===================================================================
\section{\vrev{종합 논의}}
\label{sec:ch5_discussion}
%% ===================================================================

제4장(연구 결과)의 교차분석의 분석 결과를 종합하여 DRTO 개별 시그모이드
바운딩 모델의 이론적 기여, 방법론적 함의, 그리고 한계를 논의한다.
교차분석에서 도출된 \jrev{\chg{주요 결과와 발견, 검증}은} 다음 세 가지 핵심 메시지로 통합된다.
\label{subsec:ch5_integration}

\chg{본 장의 논의는 \chref{ch:introduction}에서 제기한 다섯 가지 연구 목적에 대한 응답으로 구성된다.
서론은 관측성 단일 채널 모델을 관측성과 불확실성의 이중 채널 $\DRTO$ 모델로 확장하고,
정적 기준에서 관측성과 불확실성을 단계적으로 도입하는 네 개 운영 조건 체계(C0--C3)를 설계하며,
모델의 효과와 강건성, 회귀 설명력, 이중채널 비선형 효과, 전문가 타당성의 다섯 측면을 다층 검증하고,
관측성과 극단 불확실성이 결합할 때 나타나는 이중채널 비선형 효과를 규명하며, 6개 산업에 적용
가능한 실무 가이드라인을 제시하는 것을 목표로 제시하였다. 본 연구는 이들 목적이 제4장(연구 결과)로 어떻게
충족되는지를 5.2(학술적 기여)와 5.3(실무적 기여)로
나누어 논의하며, 제4장(연구 결과)의 모든 결과 기여와 활용 귀착은 5.1.1(제4장 연구 결과의 기여와 활용 종합 매핑)에 모두 매핑하였다. 또한 세 개의 핵심 메시지는 다음과 같다.}

\chg{첫째, 핵심 메시지는 관측성과 불확실성의 구조적 반대 효과이다.}
C0$\to$C1 전환(Cohen's $d = -1.180$)과 C1$\to$C2 전환($d = +1.871$)의
효과 크기는 모두 \jrev{\textbf{대효과(large effect)}} 수준이며 방향이 반대이다.
이는 개별 시그모이드 바운딩 모델에서 $E_{O,\text{bnd}} \leq 0$과
$E_{U,\text{bnd}} \geq 0$으로 설계된 구조적 반대 효과가
실험적으로 확인된 것이다.
\pdferr{2차 회귀}에서 C1은 $R^2 = 1.000$, C2는 $R^2 = 0.998$에 도달하여,
이 반대 효과가 \pdferr{2차 다항식}의 정밀도로 정량화됨을 보여준다.
실무적으로 이 발견은 BCM 전략 수립 시 관측성 투자(복구 시간 단축)와
불확실성 대비(안전 마진 확보) 사이의 trade-off를 명확히 인식해야
함을 시사한다.

\chg{둘째, 핵심 메시지는 평균의 비차별성과 극단의 극적 분기이다.}
C2--C3 비교에서 Cohen's $d = -0.238$(\jrev{소효과(small effect)})이지만,
$\mathrm{CVaR}_{99}$ 격차는 $+24.0\%$, Tail $k_O$ 계수 감쇠는 $0.52$배이다.
이 ``평균에서의 비차별성과 극단에서의 극적 분기'' 패턴은
불확실성의 분포 형태가 평균적 의사 결정이 아닌 극단 위험 관리에서
결정적으로 중요함을 정량적으로 입증한다.
P85.8 교차점은 \jrev{이 두 체계의 분기점(자연적 경계)}이며,
분할 회귀에서 Core($R^2 = 0.999$)와 Tail($R^2 = 0.710$)의
높은 설명력은 \jrev{이 분기의 통계적 정당성}을 확인한다.

\chg{셋째, 핵심 메시지는 전문가 검증에 의한 실무적 타당성 확인이다.}
5명의 전문가가 전체 평균 $4.46/5.0$으로 프레임워크를 평가하였으며
($t = 3.68$, $p = 0.011$), 15개 하위 주제 전체에서 $80\%$ 이상
동의가 확인되었다. 이 결과는 시뮬레이션에 기반한 이론적 발견이
BCM 실무자의 경험적 관찰과 부합함을 삼각 검증(triangulation)한 것이다.
특히 전문가들이 P85.8 교차점과 이중 채널 감쇠 메커니즘의 실무적
설명력을 높이 평가한 점은 이 발견의 이론적 신규성과 실무적
유용성을 동시에 확인한다.


\subsection{\chg{제4장 연구 결과의 기여와 활용 종합 매핑}}
\label{subsec:ch5_results_map}

\chg{본 항은 \chref{ch:results}에서 도출된 모든 결과가 학술적 기여와 실무적 활용 가운데 어디로 귀착되는지를
일람으로 제시한다. 이는 의미의 크고 작음과 무관하게 각 결과가 본 장의 어느 논의에서 성과로 활용되고
기여로 정리되는지를 명시함으로써, 결과와 논의 사이의 누락 없는 연결을 보장하기 위함이다. 학술적 기여는
\secref{sec:academic-contrib}에서, 실무적 활용은 \secref{sec:practical-contrib}에서 상세히 전개되며,
\tabref{tab:ch5-results-map}\refpo{tab:ch5-results-map}{은}{는} 그 매핑을 결과 영역별로 정리한 것이다.}

\begin{table}[htbp]
  \centering
  \caption{\chg{제4장 결과 영역별 학술적 기여와 실무적 활용 귀착}}
  \label{tab:ch5-results-map}
  \small
  \begin{tabularx}{\textwidth}{l X X}
    \toprule
    \textbf{\chg{결과 영역(4장 절)}} & \textbf{\chg{핵심 결과}} & \textbf{\chg{학술 기여 / 실무 활용 귀착}} \\
    \midrule
    \chg{관측성 효과 (4.1)} & \chg{$\bar\eta_{\text{C1}}=0.853$($-14.7\%$), $d=-1.180$, H1$\cdot$H3-a 100\%} & \chg{이론(5.2.1)$\cdot$실증(5.2.3) / 회귀식 산출(5.3.1)$\cdot$성숙도 진단(5.3.3)} \\
    \addlinespace
    \chg{가우시안 프리미엄 (4.2)} & \chg{$\bar\eta_{\text{C2}}=1.682$($+68.2\%$), $d=+1.871$, 프리미엄율 84.9\%} & \chg{이론(5.2.1)$\cdot$실증(5.2.3) / 회귀식 산출$\cdot$안전마진 해석(5.3.1)} \\
    \addlinespace
    \chg{파레토$\cdot$극단 위험 (4.3)} & \chg{$\bar\eta_{\text{C3}}=1.514$, P85.8 CDF 교차점, CVaR$_{99}$ $+24.0\%$, CVaR 에스컬레이션} & \chg{이론(5.2.1) / \textbf{이중채널 자원배분$\cdot$극단 대비(5.3.2)}} \\
    \addlinespace
    \chg{회귀$\cdot$간명모형 (4.4)} & \chg{$R^2$ 진행(C1 1.000/C2 0.998/C3 0.858), 분할 Core 0.999, 4변수 간명모형, 지배변수 전환} & \chg{방법론(5.2.2) / \textbf{회귀식 DRTO 산출 도구(5.3.1)}$\cdot$산업별 투자우선순위(5.3.4)} \\
    \addlinespace
    \chg{이중채널 감쇠 (4.10)} & \chg{Tail $k_O$ $-0.415$ vs Core $-0.798$($0.52$배), Core 선형/Tail 포화 메커니즘} & \chg{이론(5.2.1) / \textbf{이중채널 스트레스테스트$\cdot$자원배분(5.3.2)}} \\
    \addlinespace
    \chg{\m{효과 크기}$\cdot$분포구조 (4.7--4.9)} & \chg{경로별 Cohen's $d$, 순효과 격차(가우시안 1.110 vs 파레토 0.650), 백분위$\cdot$PDF 교차, 회귀계수 직접감쇠($k_O$ 0.75배$\cdot$U 0.45배)} & \chg{이론(5.2.1)$\cdot$실증(5.2.3) / 산업별 환경계수$\cdot$투자효과 비교(5.3.4)} \\
    \addlinespace
    \chg{강건성$\cdot$기반검증 (4.5--4.6)} & \chg{기술통계 질적전환, $n{=}10{,}000$ 적정성, $10^8$ 다중시드(H6$\cdot$H7), Golden Compare 47점(L0)} & \chg{방법론(5.2.2)$\cdot$실증(5.2.3) / 재현성 신뢰 근거(5.3.3)} \\
    \addlinespace
    \chg{전문가 검증 (4.11)} & \chg{$4.46/5.0$($p{=}0.011$, $d{=}1.65$), CVI $0.985$, $\alpha{=}0.89$, 인터뷰 15/15 $\geq80\%$(H9$\cdot$H10)} & \chg{실증(5.2.3) / 실무 수용성$\cdot$현장 통찰(5.3 전반)} \\
    \addlinespace
    \chg{교차 정합성 (4.12)} & \chg{\citet{lee2026drto} 매개변수 승계 일관성, 303개 교차검증 PASS, 한계-해결 연쇄} & \chg{방법론(5.2.2)$\cdot$실증(5.2.3) / 연구 신뢰성 보증} \\
    \addlinespace
    \chg{가설 종합 (4.13)} & \chg{H1--H10 전원 채택, \frev{분포 환경 적합성} 종합 명제, $\kappa{=}1.2$ 잠정 채택} & \chg{실증(5.2.3)$\cdot$한계(5.4) / $\kappa$ 기본값 적용 근거(5.3)} \\
    \bottomrule
  \end{tabularx}
\end{table}

\chg{이 매핑이 보이는 핵심은, 제4장(연구 결과)의 결과가 학술적으로는 이론, 방법론, 실증의 세 기여로, 실무적으로는
회귀식 기반 산출, 이중채널 기반 자원배분, 성숙도와 에스컬레이션 운영, 산업별 적용의 네 도구로 빠짐없이
이어진다는 점이다. 특히 본 연구의 두 핵심 발견인 \textbf{회귀식(간명모형)}과 \textbf{이중채널 감쇠
메커니즘}은 각각 5.3(실무적 기여)의 독립된 실무 도구로 구체화되며, 이는 이론적 발견이
실무 의사결정으로 내려오는 경로를 명시한다.}




\bigskip
이상의 교차분석과 논의를 종합하면, $\DRTO$ 개별 시그모이드 바운딩 모델은
관측성과 불확실성의 효과를 구조적으로 분리하고, 극단 영역에서의
비선형 상호작용을 이중채널 감쇠로 포착하며,
전문가 검증에 의해 실무적 타당성이 확인된 프레임워크이다.
이어지는 절에서는 이러한 이론적 발견이 BCM 실무와 정책 수립에 갖는
구체적 시사점을 논의한다.


%% ===================================================================
%% 이하 구 제6장(시사점) 내용 — ch5와 병합
%% ===================================================================

이하에서는 \chref{ch:results}의 시뮬레이션 실험과 앞서 논의한 교차 조건 분석 결과를
토대로 $\DRTO$ 프레임워크의 시사점을 이론적, 방법론적, 그리고 실무적 관점에서
체계적으로 논의한다. 먼저 개별 시그모이드 바운딩 모델이 BCM 이론에
기여하는 바를 정리하고(\secref{sec:theoretical-implications}),
연구 방법론의 혁신과 확장 가능성을 논의한다(\secref{sec:methodological-implications}).
이어서 2단계 프레임워크, 에스컬레이션 체계, 6개 산업 적용 가이드라인 등
실무적 시사점을 제시하고(\secref{sec:practical-implications}),
마지막으로 연구의 한계점을 투명하게 논의한다(\secref{sec:limitations}).

                   % 5.1 종합 논의 (5.1.2/5.1.3 분리됨)
% =====================================================================
% 5.2 학술적 기여 — 5단계 5-2 (clean subsection wrapper)
% 본문 텍스트 무수정. 분할 본문 파일을 새 subsection 구조로 결합.
%   5.2.1 이론적 기여  ← 5.1.2 + 5.2(전체) + 5.8.1 본문
%   5.2.2 방법론적 기여 ← 5.3(전체) + 5.8.2 본문
% =====================================================================

\section{\vrev{학술적 기여}}
\label{sec:academic-contrib}

\begin{p5add}
본 학술적 기여는 본 연구가 BCM\chg{과 재해 복구}(DR) 학문 분야의
지식 체계에 새롭게 기여한 바를 \chg{\m{(i)} 이론적, \m{(ii)} 방법론적, 그리고 \m{(iii)} 실증적} 세 차원에서
종합한다. 이론적 기여(5.2.1항)에서는 $\DRTO$ 모델의 수학적 구조와
이중 채널 감쇠 메커니즘의 신규성을, 방법론적 기여(5.2.2항)에서는 비중첩
대역 시드 체계와 분할 회귀 방법론을, 실증적 기여(5.2.3항)에서는 \jrev{10개
가설} 전원 채택과 다중 검증 프레임워크의 결과\jrev{, 그리고 \textbf{분포 환경
적합성 종합 명제}(\secref{sec:ch4_hypothesis_summary})}를 다룬다.
\end{p5add}

% --- 5.2.1 이론적 기여 ---
\subsection{\vrev{이론적 기여}}
\label{subsec:academic-theoretical}

\begingroup
  % 원본 파일의 \section, \subsection 헤더만 시각적으로 억제 (label은 유지)
  \renewcommand{\section}[1]{}
  \renewcommand{\subsection}[1]{}
  \renewcommand{\subsubsection}[1]{\paragraph{#1}}

  % --- 주축: 5.1.2 이론적 기여 (간결한 intro + 1)/2)/3) 형식) ---
  %% -------------------------------------------------------------------
\subsection{\vrev{이론적 기여}}
\label{subsec:ch5_theoretical}
%% -------------------------------------------------------------------

본 연구의 이론적 기여는 네 가지 차원에서 정리된다.

\chg{첫째, 개별 시그모이드 바운딩에 의한 물리적 경계 보장이다.}
앞서 비교한 선형 RTO 모델의 물리적 경계 위반 문제를 개별 시그모이드
바운딩 함수로 해결하였다. $E_{O,\text{bnd}} \in (-R_0, 0]$과
$E_{U,\text{bnd}} \in [0, R_{\max} - R_0)$의 범위가 해석적으로 보장되어,
$\DRTO \in (0, R_{\max})$가 입력값에 무관하게 성립한다.
이는 기존 BCM 문헌에서 제시되지 않았던 수학적 보장으로서,
\chg{복구목표시간}의 음수 산출이나 MTPD 초과라는 구조적 문제를
근원적으로 해결한다.

\chg{둘째, 이중채널 감쇠 효과의 발견과 정량화이다.}
P85.8 교차점을 기준으로 $k_O$ 계수가 $0.52$배로 감쇠되는
이중채널 현상의 발견은 BCM 이론에 새로운 통찰을 제공한다.
이 발견은 극단 불확실성과 관측성이 비선형적으로 상호 작용하는
메커니즘을 규명하며, 시그모이드 비선형성에 의한 이론적 설명을
제시한다.

\chg{셋째, \citet{lee2026drto}과 본 연구의 일관적 이론 체계 구축이다.}
선행 연구\citep{lee2026drto}의 선형 근사$\to$시그모이드 바운딩$\to$$\kappa$ 최적화에서
본 연구의 이중채널 확장(C2, C3)까지의 단계적 확장이
\chg{수식·데이터·결론}의 교차 정합성을 유지하며 수행되었다.
선형 모델이 시그모이드 바운딩의 1차 근사임을 보이는 관계
($|z| \ll 1$에서 $S(z) \approx z$)는 이론적 확장의
내적 일관성을 수학적으로 보장한다.

  넷째, 불확실성의 확률적 정량화이다. 복구 과정의 불확실성을 일상적 변동을 나타내는 가우시안과 극단 사건을 나타내는 파레토의 두 분포로 정량화함으로써, 관측성 단일 차원에 머물던 기존 모델의 한계를 보완하고 극단 영역의 꼬리 위험을 복구목표시간 산정에 반영하였다.



  % --- 보완: 5.2 이론적 시사점 (기술적 세부 설명 보완) ---
  \section{\vrev{이론적 시사점}}
\label{sec:theoretical-implications}
%% ===================================================================

본 연구의 이론적 기여는 BCM 분야에 \jrev{\chg{i) 관측성 개념의 통합, ii) 불확실성의 확률적 정량화(가우시안과 파레토), iii) 개별 시그모이드 바운딩 이론의 정립, iv) 이중 채널 감쇠 효과의 발견이라는} 네 가지 축}으로 정리된다.
\p{[}표~\vref{tab:ch6-theoretical-contributions}\p{]}는 이론적 기여의 핵심 내용을 종합한다.

\begin{table}[htbp]
\centering
\caption{\jrev{이론적 기여 종합}}
\label{tab:ch6-theoretical-contributions}
\small
\begin{tabularx}{\textwidth}{c l X c}
\toprule
\textbf{\#} & \textbf{기여 영역} & \textbf{핵심 내용} & \textbf{근거} \\
\midrule
1 & 관측성 통합 &
  제어 이론의 관측성 개념을 BCM에 최초 도입;
  $z_O = \kappa \cdot k_O \cdot O_{\mathrm{raw}}$으로 정량화하여
  정적 RTO의 구조적 한계를 극복 &
  H1, \jrev{H3}-a \\
\addlinespace
2 & 시그모이드 바운딩 이론 &
  수정 시그모이드 $S(z) = 2/(1+e^{-z})-1$을 기저 함수로 채택;
  $\DRTO = R_0 + E_{O,bnd} + E_{U,bnd}$로 물리적 경계
  $(0,\; R_{\max})$를 해석적으로 보장 &
  \jrev{H2 ($\kappa = 1.2$ 전제)} \\
\addlinespace
3 & 관측성--불확실성 이원 구조 &
  관측성 효과($E_{O,bnd} \leq 0$)와 불확실성 효과($E_{U,bnd} \geq 0$)의
  구조적 반대 방향 확인; Cohen's $d$: $-1.180$ vs $+1.871$ &
  \jrev{H1, H8} \\
\addlinespace
4 & P85.8 CDF 교차점 규명 &
  가우시안(C2)과 파레토(C3) 분포의 $\eta$ 분위수 함수가
  P85.8에서 교차; 꼬리 효과의 조건부 특성 규명 &
  \jrev{H3}-c \\
\addlinespace
5 & 이중채널 감쇠 발견 &
  P85.8 기준 분할 회귀에서 Tail 영역의 $k_O$ 계수($-0.415$)가
  Core($-0.798$) 대비 $0.52$배 감쇠; 극단 환경에서 관측성
  효과의 상대적 축소 &
  \jrev{H8} \\
\addlinespace
6 & 꼬리 리스크 비선형 증폭 &
  CVaR$_{99}$: C3가 C2 대비 $+24.0\%$;
  분위수 상승에 따른 비선형 격차 확대
  (CVaR$_{90}$: $+20.4\%$ $\to$ CVaR$_{99}$: $+24.0\%$) &
  \jrev{H8} \\
\bottomrule
\end{tabularx}
\end{table}


%% -------------------------------------------------------------------
\subsection{\vrev{\revdel{BCM 이론에 대한 관측성 통합}}}
\label{subsec:observability-integration}
%% -------------------------------------------------------------------

\fdel{전통적 BCM 이론에서 RTO는 BIA(업무영향분석) 시점에 고정되는
정적 매개변수로 취급되어, 시스템 상태의 실시간 변화를 반영하지 못하였다.
본 연구는 제어 이론~\citep{kalman1960control}의 관측성(observability) 개념을
BCM 분야에 최초로 도입하여, 시스템의 내부 상태를 외부 출력으로부터
추정할 수 있는 정도를 $O_{\mathrm{raw}} \in [0,1]$로 정규화하고,
이를 관측성 가중치 $k_O$와 곡률 매개변수 $\kappa = 1.2$를 통해
$\DRTO$의 관측성 조정량으로 변환하였다.}

\fdel{관측성 조정량 $E_{O,bnd} = -R_0 \cdot S(\kappa \cdot k_O \cdot O_{\mathrm{raw}})$는
항상 비양수($\leq 0$)이므로, 관측성이 높을수록 $\DRTO$가 $R_0$ 이하로 단축된다.
이 구조는 \jrev{H3}-a에서 $100\%$ 준수율로 검증되었으며, 이는 확률적 결과가 아닌
모델의 \textbf{구조적 필연}이다. Cohen's $d = -1.180$의 매우 큰 효과 크기는
관측성 통합이 BCM의 복구 예측 정확도를 질적으로 향상시킨다는 이론적 근거를 제공한다.}

\fdel{이러한 관측성 통합은 기존 BCM 이론의 세 가지 공백을 해소한다.
첫째, RTO를 ``고정 상수''에서 ``상태 의존 함수''로 전환함으로써
동적 환경에서의 BCM 의사결정 기반을 확보한다.
둘째, APM(Application Performance Monitoring), SIEM(Security Information and
Event Management) 등 기존 모니터링 인프라의 BCM 활용 경로를 이론적으로 정당화한다.
이는 SRE(Site Reliability Engineering) 분야에서 관측성을 에러 예산(error budget)과
연계하여 서비스 신뢰성을 정량적으로 관리하는 접근(Beyer et al.\ 2016)과
맥락을 공유하며, DevOps 문화에서 모니터링 인프라 투자가 복구 성능을 향상시킨다는
실증적 발견(Forsgren et al.\ 2018; Kim et al.\ 2016)과도 부합한다.
셋째, 관측성 투자($k_O$ 증가)와 불확실성 허용($k_U$ 증가) 간의
trade-off를 $k_O + k_U = 1$ 제약으로 형식화하여,
BCM 자원 배분의 분석적 프레임워크를 제시한다.}


%% -------------------------------------------------------------------
\subsection{\vrev{\revdel{개별 시그모이드 바운딩 이론}}}
\label{subsec:sigmoid-bounding-theory}
%% -------------------------------------------------------------------

\fdel{$\DRTO$ 모델의 수학적 핵심은 수정 시그모이드 함수
$S(z) = 2/(1+e^{-z})-1 \in (-1,1)$에 의한 개별 바운딩(individual bounding)이다.
관측성 효과와 불확실성 효과를 각각 독립적인 시그모이드로 바운딩하여
$\DRTO = R_0 + E_{O,bnd} + E_{U,bnd}$로 정식화함으로써, 다음의 세 가지
이론적 요구사항을 동시에 충족한다.}

% [5단계 Plan H 재편: 구 5.2.2 "개별 시그모이드 바운딩 이론" 3개 항목
%  (물리적 경계 보장 / 무한 미분 가능성 / 교호작용 포착)은 5.2.1 이론적
%  기여로 통합·흡수됨. 기존 \fdel{} 래핑된 본문은 빈 enumerate 번호(1/2/3)
%  잔재를 발생시켜 enumerate 블록 전체 제거.]



%% -------------------------------------------------------------------
\subsection{\vrev{\revdel{이중채널 감쇠 효과와 P85.8 교차점}}}
\label{subsec:dual-channel-crossover}
%% -------------------------------------------------------------------

\fdel{본 연구의 가장 핵심적인 이론적 발견은 P85.8 분할 기준점(CDF 교차점과 일치)을
경계로 하는 이중채널 감쇠(dual-channel attenuation) 효과이다.
C3 조건에서 P85.8 기준으로 분할 회귀를 수행한 결과,
관측성 가중치 $k_O$의 회귀 계수가 Core 영역($\leq$ P85.8)의 $-0.798$에서
Tail 영역($>$ P85.8)의 $-0.415$로 $0.52$배 감쇠된다.}

\fdel{이 발견의 이론적 함의는 세 가지이다.
첫째, 극단적 불확실성 환경에서 관측성 투자의 \textbf{한계효용이 감소}한다.
Core 영역에서 $k_O$를 $0.1$ 증가시키면 $\eta$가 $0.080$ 감소하나,
Tail 영역에서는 동일한 변화가 $\eta$를 $0.042$만 감소시킨다.
이는 BCM 자원 배분에서 극단 시나리오일수록 관측성 투자의
비용-효과성이 감소하며 불확실성 관리가 우선되어야 함을 의미한다.
둘째, 시그모이드 비선형성과 파레토의 \textbf{상호작용 메커니즘}이
규명되었다. Tail 영역에서 큰 $U$ 값이 시그모이드의 포화 구간을 활성화하여
불확실성 채널이 지배적이 되며, $k_O$의 상대적 역할이 축소된다.
셋째, P85.8이 분할 회귀 기준점이자 CDF 교차점으로서 가우시안--파레토 분포 전환의
\textbf{구조적 경계}임이 다중 시드 검증을 통해 확인되어,
이 교차점이 특정 시드에 의존하지 않는 강건한 현상임이 입증되었다.}


  % --- 5.8.1 (이론적 기여) — 5.1.2와 3중 주제 중복으로 5-2 삭제 ---
  \begin{p5del}{현 5.8.1 이론적 기여 — 5.1.2의 1)/2)/3)과 동일 주제 중복}
  %% ===================================================================
\section{\vrev{학술적 기여}}
\label{sec:conclusion-contributions}
%% ===================================================================


%% -------------------------------------------------------------------
\subsection{\vrev{이론적 기여}}
\label{subsec:contrib-theoretical}
%% -------------------------------------------------------------------

\fdel{본 연구의 이론적 기여는 다음의 세 가지로 정리된다.}

\fdel{첫째, \textbf{업무연속성관리$\cdot$관측성$\cdot$시그모이드 바운딩의 통합적 이론 기반을
확립}하였다.
기존 BCM 연구에서 RTO는 BIA 시점에 고정되는 정적 매개변수로 취급되었으며,
관측성과 불확실성을 단일 수리 모델에서 통합한 선행 연구는 부재하였다.
본 연구는 제어 이론의 관측성 개념과 확률론적 불확실성 모형을
수정 시그모이드 함수라는 공통 기저 함수 위에서 결합함으로써,
BCM 분야에 새로운 이론적 기반을 제공한다.
특히 $\DRTO = R_0 + E_{O,\text{bnd}} + E_{U,\text{bnd}}$의 가산 구조는
각 효과의 독립적 분석과 결합 효과의 동시 포착을 가능하게 하여,
모델의 해석 가능성(interpretability)과 확장 가능성(extensibility)을
동시에 확보한다.}

\fdel{둘째, \textbf{꼬리 감쇠 현상을 발견하고 정량화}하였다.
P85.8 분할 기준을 경계로 관측성 가중치 $k_O$의 효과가
Core($-0.798$)에서 Tail($-0.415$)로 $0.52$배 감쇠되는 현상은,
극단 불확실성 환경에서 시그모이드 포화에 의해
관측성 채널의 조절 효과가 구조적으로 제한됨을 규명한 것이다.
이 발견은 BCM 자원 배분 전략에 새로운 분석적 기반을 제시하며,
극단 상황에서의 대비 전략이 관측성 투자만으로는 한계가 있으며
사전 자원 확보가 병행되어야 한다는 실무적 시사점을 제공한다.}

\fdel{셋째, \textbf{P85.8 CDF 교차점을 가우시안--파레토 분포 전환의 구조적 경계로 규명}하였다.
C2와 C3의 CDF가 P85.8($\eta \approx 2.42$)에서 교차하며,
이 교차점을 경계로 두 조건의 행태가 질적으로 전환된다.
P85.8 이하에서 두 분포는 사실상 구분 불가능하나,
P85.8 초과에서 C3의 $\eta$가 C2를 급격히 상회하여
CVaR$_{99}$ 기준 $+24.0\%$의 격차를 보인다.
이는 파레토 불확실성의 효과가 \textit{평균 수준의 비차별성}과
\textit{꼬리 영역의 극적 분기}로 구성되는 이중 구조를 가진다는 것을
실증적으로 규명한 결과이다.}


  \end{p5del}
\endgroup

% --- 5.2.2 방법론적 기여 ---
\subsection{\vrev{방법론적 기여}}
\label{subsec:academic-methodological}

\begingroup
  \renewcommand{\section}[1]{}
  \renewcommand{\subsection}[1]{}
  \renewcommand{\subsubsection}[1]{\paragraph{#1}}

  

%% ===================================================================
\section{\vrev{방법론적 시사점}}
\label{sec:methodological-implications}
%% ===================================================================

본 연구는 $\DRTO$ 프레임워크의 검증 과정에서 \jrev{네 가지} 방법론적 혁신을
달성하였다. \chg{이를 차례로 서술한다.}


%% -------------------------------------------------------------------
\subsection{\vrev{몬테카를로 시뮬레이션 방법론}}
\label{subsec:mc-methodology}
%% -------------------------------------------------------------------

\chg{첫째, 몬테카를로 시뮬레이션 방법론이다.} $n = 10{,}000$ 표본 $\times$ 4개 조건(C0--C3) $\times$ $10{,}000$개 마스터 시드의
총 $4 \times 10^8$회 $\DRTO$ 계산을 통해, 시뮬레이션 기반 BCM 모델 검증의
새로운 표준을 제시하였다.

\textbf{비중첩 대역 시드(non-overlapping band seeding)} 체계는 마스터 시드
$s = 42$로부터 변수별 시드를 $10{,}000$ 간격의 비중첩 대역\chg{(}$O_{\mathrm{raw}}$:
$[0\text{--}9999]$, $k_O$: $[20000\text{--}29999]$, $U^{(G)}$:
$[30000\text{--}39999]$, $U^{(P)}$: $[40000\text{--}49999]$\chg{)}으로
파생한다. \jrev{H6} 검증에서 6개 변수쌍 모두 $|r| < 0.02$(최대 $|r| = 0.016$)를
충족하여, 시뮬레이션의 완전한 재현성과 변수 간 통계적 독립성이 동시에 보장되었다.

이 체계의 방법론적 의의는 세 가지이다.
\chg{첫째,} 단일 시드로부터 다수의 독립적 난수열을 파생하는 표준화된 규약을 제공한다.
\chg{둘째,} $10{,}000$개 시드에 걸친 $10^8$회 다중 시드 검증(\jrev{H7})에서 Bonferroni 보정
$z$-검정 결과 45개 계수 중 기각이 $0$개로, 특정 시드에 대한 의존성을 완전히
배제한다. \chg{셋째,} 재현성 프로토콜이 명시되어 있어 후속 연구자의 독립적 검증이 가능하다.


%% -------------------------------------------------------------------
\subsection{\jrev{\texorpdfstring{\citet{lee2026drto}}{Lee and Cheung (2026)} $\kappa$ 매개변수의 본 연구 환경 적용 및 \frev{분포 환경 적합성} 시사 방법론}}
\label{subsec:kappa-optimization}
%% -------------------------------------------------------------------

\chg{둘째, $\kappa$ 매개변수의 분포 환경 적합성 시사 방법론이다.} \jrev{본 연구는 곡률 매개변수 $\kappa^{*} = 1.2$를 \chg{\citet{lee2026drto}}에서 다기준 종합점수($f_1$ 실무 유의성,
$f_2$ 비포화율, $f_3$ 효과 크기, $f_4$ 설명력)로 \chg{도출하여 출판된} 매개변수로
받아들이고, 본 연구의 새로운 분포 환경(\chg{C2 가우시안과 C3 파레토})에 그대로 적용하였다. 따라서 $\kappa^{*}$ 도출 자체는 본 연구의
독자적 기여가 아니며, \chg{\citet{lee2026drto}의} 격자 탐색 결과(C1 조건 기준
$F(\kappa^{*}) = 0.9997$)를 \textbf{전제 조건}으로 받아들인다.}

\jrev{본 연구의 방법론적 보완은 다음과 같다. \chg{\citet{lee2026drto}의}
4기준 종합점수 $F(\kappa) = f_1 \cdot f_2 \cdot f_3 \cdot f_4$는
C1 환경 기반으로 정의되어, 새로운 분포 환경(\chg{C2 가우시안과 C3 파레토})에서는
직접 격자 탐색에 부적합하다는 학술적 한계를 지닌다. 본 연구는 격자
탐색의 직접적 대안으로서, 9개 가설(\jrev{H2--H10}) PASS의 \textbf{통합 해석을
통한 \jrev{\frev{분포 환경 적합성}} 종합 명제 도출 방법론}을 제시하였다
(\secref{sec:ch4_hypothesis_summary} \jrev{\frev{분포 환경 적합성}} 종합 명제 참조).
구체적으로, \chg{i) 물리적 안전성(H2), ii) 이론적 일치(H3), iii) 통계적 설명력(H4·H5),
iv) 결과 안정성(H6·H7), v) 메커니즘 발현(H8), vi) 외부 검증(H9·H10)}의 6가지 차원에서 $\kappa = 1.2$의 \jrev{\frev{분포 환경 적합성}}이 집합적으로
\frev{시사된다}.}

\jrev{이러한 \textbf{다중 가설 PASS의 통합 해석} 방법론은 \chg{\citet{lee2026drto}}
매개변수의 새로운 환경 적용성을 \frev{시사}하는 보완적
방법론적 틀로서, 본 연구가 독자적으로 제시하는 기여이다. 본 접근법은
선행연구를 기반으로 분포 환경을 확장하는 후속 연구 일반에 적용 가능하며,
직접적 격자 탐색이 부적합한 환경에서의 매개변수 적합성 검증에 활용될 수 있다.}


%% -------------------------------------------------------------------
\subsection{\vrev{분할 회귀(Split Regression) 방법론}}
\label{subsec:split-regression}
%% -------------------------------------------------------------------

\chg{셋째,} P85.8 기준 분할 회귀는 파레토 분포의 체제 전환(regime change)을
포착하기 위한 방법론적 혁신이다. C3 전체 표본에서 \pdferr{2차 회귀}의
$R^2 = 0.858$에 그치는 잔여 비설명 분산($14.2\%$)을,
P85.8 기준으로 Core($\leq$ P85.8, $n = 8{,}580$)와 Tail($>$ P85.8, $n = 1{,}420$)로
분할함으로써 Core $R^2 = 0.999$, Tail $R^2 = 0.710$으로 개선하였다.

분할 회귀의 방법론적 기여는 다음과 같다.
\chg{i)} 단일 회귀 모형으로 포착할 수 없는 \textbf{구간별 이질적 구조}를
데이터 기반으로 식별하는 체계적 절차를 제시한다.
\chg{ii)} 분할 기준점(P85.8)의 결정이 사전적 이론(CDF 교차)에 근거하므로,
임의적 구간 분할의 자의성 문제를 해소한다.
\chg{iii)} 이 방법론은 BCM 분야에 국한되지 않고, 파레토 분포가 관여하는
금융 리스크 관리, 보험 계리, 자연 재해 모형 등에
범용적으로 확장 가능하다.

  
%% -------------------------------------------------------------------
\subsection{\vrev{방법론적 기여}}
\label{subsec:contrib-methodological}
%% -------------------------------------------------------------------

\fdel{방법론적 관점에서 본 연구는 다음의 세 가지 기여를 달성하였다.}

\fdel{첫째, \textbf{P85.8 기준 분할 회귀(split regression)를 파레토 분포 분석에 적용}하였다.
C3 전체 표본에서 2차교호작용 회귀의 $R^2 = 0.858$에 그치는 잔여 비설명 분산을,
P85.8 기준으로 Core($\leq$ P85.8)와 Tail($>$ P85.8)로 분할함으로써
Core $R^2 = 0.999$(0.9990), Tail $R^2 = 0.710$(0.7100)으로 Core 영역의
결정론적 구조를 완전히 포착하였다.
이 분할 전략은 파레토 분포의 체제 전환(regime change)을 포착하는
범용적 분석 방법론으로, BCM을 넘어 보험, 금융, 재난관리 등
파레토 현상이 핵심인 분야에 확장 적용될 수 있다.}

\fdel{둘째, \textbf{비중첩 대역 시드(non-overlapping band seeding) 체계}를 설계하였다.
마스터 시드 $s = 42$로부터 변수별 시드를 $10{,}000$ 간격의 비중첩 대역($O_{\mathrm{raw}}$: $[0\text{--}9999]$, $k_O$: $[20000\text{--}29999]$,
$U^{(G)}$: $[30000\text{--}39999]$,
$U^{(P)}$: $[40000\text{--}49999]$)으로 파생하여,
시뮬레이션의 완전한 재현성을 보장하면서도 변수 간 통계적 독립성을
유지한다(모든 변수 쌍 $|r| < 0.02$).}

\jrev{넷째}, \textbf{6단계 다층 검증(six-stage multi-layer validation) 프레임워크}를 구축하였다.
기반 검증(L0, Golden Compare), 시뮬레이션 검증, 회귀 검증, 강건성 검증, 이중채널 검증, 전문가 평가의
6개 단계(L0--L5)가 계층적 의존 구조를 형성하여 모델의 타당성을 다각적으로 보장한다.
정량적 시뮬레이션과 정성적 전문가 평가를 삼각검증(triangulation)함으로써
단일 방법론의 편향을 최소화하였으며, 이 프레임워크는
수리 모델 기반 연구의 검증 설계 표준으로 활용될 수 있다.



\endgroup

% --- 5.2.3 실증적 기여 (5-2단계 신규 작성) ---
\subsection{\vrev{실증적 기여}}
\label{subsec:academic-empirical}

\begin{p5add}
본 연구에서 실증적 기여는 시뮬레이션과 전문가 평가의 결합 검증을 통해 달성한
결과를 토대로 정리한다.

첫째, \textbf{\jrev{10개 연구 가설}(H1--\jrev{H10})이 사전 설정 기준을 모두 충족하여 모두
채택}되었다. 시뮬레이션 검증(H1--\jrev{H3}), 회귀 검증(\jrev{H4}--\jrev{H5}), 강건성 검증(\jrev{H6}--\jrev{H7}),
이중 채널 검증(\jrev{H8}), 전문가 평가(\jrev{H9}--\jrev{H10})의 6단계 다층 검증 프레임워크에서
계층적 의존 구조가 모두 만족됨으로써, $\DRTO$ 모델의 내부 정합성과
외부 타당성이 동시에 확인되었다.

둘째, \textbf{$10{,}000$개 마스터 시드에 걸친 $10^8$회 다중 시드 검증}을
통해 결과의 시드 의존성을 완전히 배제하였다. Bonferroni 보정 $z$-검정에서
45개 회귀 계수 중 기각이 0개로, 시뮬레이션 결과가 특정 난수열에 의존하지 않는
강건성을 정량적으로 입증하였다.

셋째, \textbf{BCM/DR 분야 전문가 5명에 의한 \chg{정량적 및 정성적} 혼합 검증}을 통해
실무 타당성을 확보하였다. 27문항 5점 리커트 척도에서 전체 평균
$4.46/5.0$($t(4) = 3.68$, $p = 0.011$, Cohen's $d = 1.65$)을 기록하였으며,
S-CVI/Ave $= 0.985$, Cronbach's $\alpha = 0.89$로 내용 타당도와 내적
일관성이 확인되었다. 심층 인터뷰에서는 5대 주제 15개 하위 주제 모두에서
$\geq 80\%$ 동의가 달성되어, 시뮬레이션 결과가 실무자의 경험적 관찰과
부합함을 삼각 검증하였다.

\jrev{넷째, \textbf{\frev{분포 환경 적합성} 종합 명제}를 도출하였다. \chg{\citet{lee2026drto}에서 격자
탐색으로 도출하여 출판한} $\kappa^{*} = 1.2$를 본 연구의
새로운 분포 환경에 적용한 결과,
9개 가설(\jrev{H2--H10})의 PASS가 \chg{i)} 물리적 안전성, \chg{ii)} 이론적 일치, \chg{iii)} 통계적
설명력, \chg{iv)} 결과 안정성, \chg{v)} 메커니즘 발현, \chg{vi)} 외부 검증의 6가지 측면에서
$\kappa = 1.2$의 \jrev{\frev{분포 환경 적합성}}을 집합적으로 \frev{시사한다} (\secref{sec:ch4_hypothesis_summary} 참조).
신규 격자 탐색을 직접 수행하지 않고도 다중 가설 PASS의 통합 해석을 통해
\chg{\citet{lee2026drto}의 매개변수의} 본 연구 환경 적용성을 학술적으로 정직하게 \frev{시사한 것이}
본 연구의 \chg{방법론적 및 실증적} 차별점이다.}
\end{p5add}
       % 5.2 학술적 기여
% =====================================================================
% 5.3 실무적 기여 — 3차심사 대규모 재설계 (발견→도구 축)
%   5.3.1 회귀식 기반 DRTO 산출·예측 도구          [신설]
%   5.3.2 이중채널 감쇠 기반 자원배분·스트레스테스트 [신설]
%   5.3.3 η 기반 평상시 성숙도·위기시 에스컬레이션   [구 5.3.1 성숙도 + 구 5.3.3 에스컬레이션 통합]
%   5.3.4 산업별 적용 가이드라인                    [구 5.3.2]
% =====================================================================

\section{\vrev{실무적 기여}}
\label{sec:practical-contrib}
\label{sec:practical-implications}    % 구 라벨 호환 — 외부 cross-ref 유지

\begin{p5add}
\chg{본 절은 본 연구가 BCM 실무 현장에 새롭게 제공한 의사결정 도구를 네 항으로 정리한다. 핵심은 본 연구의
두 핵심 발견인 회귀식(간명 모형)과 이중채널 감쇠 메커니즘을 각각 독립된 실무 도구로 구체화하여, 이론적
발견이 현장의 산출과 자원 배분으로 직접 내려오는 경로를 제시하는 데 있다.}
\chg{먼저 회귀식 기반 $\DRTO$ 산출과 예측 도구(5.3.1항)에서는 간명 회귀 모형으로 조직이 자기 환경의 $\DRTO$를
직접 계산하는 절차를, 이중채널 감쇠 기반 자원배분과 스트레스 테스트(5.3.2항)에서는 P85.8 분할과 $0.52$배
감쇠를 자원 배분 규칙과 극단 시나리오 설계로 전환하는 방안을 제시한다. 이어 $\eta$ 기반 평상시 성숙도와
위기시 에스컬레이션(5.3.3항)에서는 단일 지표 $\eta$의 평상시 진단과 위기시 대응의 이중 활용 및 선순환
구조를, 산업별 적용 가이드라인(5.3.4항)에서는 6개 대표 산업의 권장 파라미터와 분포 환경별 적용 원칙을
제시한다. 이들은 모두 \chref{ch:results}의 시뮬레이션과 회귀분석, 전문가 검증에서 도출된 정량적 결과를
실무 운용 형태로 변환한 산출물이다.}
\end{p5add}


% ===================================================================
% 5.3.1 회귀식 기반 DRTO 산출·예측 도구 [신설]
% ===================================================================
\subsection{\chg{회귀식 기반 $\DRTO$ 산출과 예측 도구}}
\label{subsec:practical-regression-tool}

\chg{본 항은 \chref{ch:results}의 회귀분석에서 도출된 간명 회귀 모형을 실무 현장의 $\DRTO$ 산출과 예측 도구로
제시한다. 이는 본 연구의 핵심 산출물인 회귀식이 단순한 통계적 적합을 넘어, 조직이 자기 환경의 $\DRTO$를
직접 계산하는 운영 도구로 활용되는 경로를 구체화한 것이다.}

\chg{도구의 토대는 전진 선택법과 LASSO의 교차검증 합의로 도출된 4변수 간명 모형(식~(\vref{eq:c2_parsimonious}),
(\vref{eq:c3_parsimonious}))이다. 9변수 완전 2차 모형이 C2에서 $R^2 = 0.998$, C3에서 $R^2 = 0.858$의
설명력을 보이는 데 비해, 합의 4변수 간명 모형은 C2 $R^2 = 0.996$, C3 $R^2 = 0.839$로 설명력 손실을
최소화하면서 해석 가능성을 확보한다. 특히 C1과 C2의 2차교호작용 회귀가 각각 $R^2 = 1.000$, $0.998$에
도달한다는 결과는 개별 시그모이드 바운딩 함수가 다항식으로 거의 완전히 근사됨을 의미하며, 이는 회귀식
산출 결과의 신뢰성을 뒷받침한다.}

\chg{실무 산출 절차는 다음과 같다. 조직은 자기 환경의 네 입력값, 곧 기준선 복구시간 $R_0$, 관측성 원시값
$O$, 관측성 가중치 $k_O$, 불확실성 수준 $U$를 간명 회귀식에 대입하여 효율비 추정값 $\hat\eta$를 얻고,
$\DRTO = \hat\eta \cdot R_0$로 동적 복구목표시간을 즉시 산출한다. 이렇게 산출된 $\DRTO$를 기존 BIA에서
고정값으로 설정한 정적 RTO와 비교하면, 관측성 투자와 불확실성 환경이 복구 목표에 미치는 영향을 정량적으로
확인할 수 있다. \tabref{tab:ch6-drto-examples}\refpo{tab:ch6-drto-examples}{은}{는} 뒤에 나올 산업별 권장
파라미터(\tabref{tab:ch6-industry-params})를 회귀식에 대입한 구체적 산출 예시이다.}

\begin{table}[htbp]
\centering
\caption{\chg{산업별 대표 $\DRTO$ 산출 예시}}
\label{tab:ch6-drto-examples}
\small
\begin{tabular}{l c c c c c}
\toprule
\textbf{산업} & \textbf{$R_0$ (h)} & \textbf{$O$} & \textbf{$\hat{\eta}$} & \textbf{$\DRTO$ (h)} & \textbf{조건} \\
\midrule
제조업   & 4.0 & 0.60 & $1.664$ & $6.66$ & C2 \\
금융     & 1.0 & 0.78 & $1.569$ & $1.57$ & C3 \\
의료     & 2.0 & 0.68 & $1.569$ & $3.14$ & C3 \\
IT 서비스 & 0.5 & 0.83 & $1.599$ & $0.80$ & C2 \\
물류     & 6.0 & 0.53 & $1.684$ & $10.10$ & C2 \\
소매     & 3.0 & 0.58 & $1.669$ & $5.01$ & C2 \\
\bottomrule
\end{tabular}
\end{table}

\chg{산출에는 4변수 간명 모형(C2 $R^2 = 0.996$, C3 $R^2 = 0.839$)을 적용하였으며, $U$는 산업별 불확실성
대표값(중앙값), $O_{\mathrm{raw}} = 0.5$를 가정하였다. 예컨대 금융 산업의 $\DRTO = 1.57$h는 $R_0 = 1.0$h
대비 $57\%$의 안전 마진을 포함한 것으로 파레토 불확실성에 의한 꼬리 위험을 반영한 값이며, 물류 산업의
$\DRTO = 10.10$h는 $R_0 = 6.0$h 대비 $68\%$ 증가한 것으로 비교적 낮은 관측성($O = 0.53$)과 불확실성
프리미엄이 결합된 결과이다. 이처럼 실무자는 자기 조직의 파라미터를 대입하여 $\DRTO$를 즉시 산출하고
기존 정적 RTO와 비교함으로써 동적 복구목표시간의 실무적 유용성을 직관적으로 확인할 수 있다.}

\chg{나아가 회귀식의 변수 진입 순서는 산출을 넘어 투자 우선순위 신호로도 활용된다. 4장의 전진 선택법
결과, 가우시안 환경(C2)에서는 관측성 가중치 $k_O$가 1순위로 진입하여 관측성이 $\eta$ 변동의 주된 결정
요인인 반면, 파레토 환경(C3)에서는 불확실성 $U$가 1순위로 진입하여 지배 변수가 전환된다. 따라서 조직은
자기 환경의 분포 특성에 따라, 가우시안 환경에서는 관측성 투자를, 파레토 환경에서는 불확실성 관리를
우선해야 한다는 자원 배분 신호를 회귀식으로부터 직접 읽어낼 수 있다.}

\chg{다만 간명 회귀식은 평균 영역의 거동을 근사하는 모형이므로, 극단 불확실성이 지배하는 꼬리 영역에서는
산출값의 정밀도가 제한된다. 꼬리 영역에서의 비선형 감쇠와 그에 따른 자원 배분은 다음 항(\subsecref{subsec:practical-dualchannel-tool})의
이중채널 도구에서 다룬다.}


% ===================================================================
% 5.3.2 이중채널 감쇠 기반 자원배분·스트레스테스트 [신설]
% ===================================================================
\subsection{\chg{이중채널 감쇠 기반 자원배분과 스트레스 테스트}}
\label{subsec:practical-dualchannel-tool}

\chg{본 항은 4장에서 발견한 이중채널 감쇠 메커니즘을 실무의 자원 배분 의사결정과 스트레스 테스트 설계
도구로 전환한다. 이중채널 감쇠란 P85.8 분할점을 경계로 관측성 가중치 $k_O$의 회귀 계수가 핵심 영역(Core,
$-0.798$)에서 꼬리 영역(Tail, $-0.415$)으로 $0.52$배 감쇠하는 현상으로, 극단 불확실성 환경에서 시그모이드
포화에 의해 관측성 투자의 한계 효용이 줄어든다는 것을 정량화한 본 연구의 핵심 발견이다.}

\chg{첫째, P85.8 분할점은 조직이 자기 불확실성 환경이 어느 체제에 있는지를 진단하고 대응 모드를 전환하는
기준이 된다. 가우시안(C2)과 파레토(C3) 분포의 누적분포함수가 P85.8($\eta \approx 2.42$)에서 교차하므로,
이 지점 이하의 일상 운영(Core)에서는 두 환경이 사실상 구별되지 않으나(Cohen's $d \approx 0$), 이 지점을
넘어서면 파레토 꼬리가 급격히 발현하여 $\mathrm{CVaR}_{99}$ 기준 C3가 C2를 $+24.0\%$ 상회한다. 따라서
조직의 복구 시간 분포가 P85.8 분할점을 넘어서기 시작하는 시점, 곧 가우시안 환경에서 파레토 환경으로의
전환은 그 자체로 자원 배분을 재조정해야 할 신호로 취급된다.}

\chg{둘째, 이중채널 감쇠 비율은 영역별 자원 배분 규칙의 정량적 근거를 제공한다. Core 영역에서는 관측성
가중치를 $0.1$ 높이면 $\eta$가 $0.080$ 감소하여 관측성 투자의 한계 수익이 충분히 발현되므로, 일상 운영의
복구 역량 향상은 관측성 투자(모니터링 확대, APM/SIEM 연동)에 집중하는 것이 효율적이다. 반면 Tail 영역에서는
동일한 관측성 개선이 $\eta$를 $0.042$만 감소시켜 한계 수익이 Core의 $0.52$배로 감쇠하므로, 극단 상황에서는
관측성 투자만으로 복구 시간을 충분히 단축하기 어렵다. 이 영역에서는 불확실성 자체의 관리, 곧
업무연속성계획(BCP) 시나리오의 다양화, 대체 사이트 확보, 사이버 보험과 같은 꼬리 위험 전가 수단에 자원을
전환해야 한다. 이는 동일한 예산이라도 환경 체제에 따라 투자 대상이 질적으로 달라져야 함을 의미한다.}

\chg{셋째, 이 발견은 스트레스 테스트 설계의 기준을 제시한다. $0.52$배 감쇠는 평균 영역의 시나리오만으로는
극단 환경에서의 복구 역량을 검증할 수 없음을 뜻하므로, BCM 스트레스 테스트는 Tail 영역의 시나리오를 별도로
설계해야 한다. 구체적으로, 파레토 분포의 P99 수준 극단값을 입력으로 하는 시나리오에서 관측성 투자만으로는
목표 복구 시간을 달성하지 못함을 검증하고, 불확실성 관리 수단을 병행했을 때의 효과를 함께 측정하는 이중
시나리오 설계가 요구된다.}

\chg{넷째, 환경 전환에 따른 회귀 계수 변화는 모델 재교정의 신호로 활용된다. 4장의 회귀계수 직접 비교 결과,
가우시안에서 파레토로 환경이 전환되면 관측성 변수 $k_O$의 계수가 $0.75$배, 불확실성 변수 $U$의 선형 계수가
$0.45$배로 변화한다. 또한 관측성 투자의 순효과는 가우시안 경로(Cohen's $d = +1.110$)보다 파레토 경로(Cohen's
$d = +0.650$)에서 상대적으로 작게 나타나, 환경에 따라 관측성 투자의 순효과 크기가 달라진다는 점도 확인된다.
조직은 자기 환경의 분포가 전환되는 것을 감지하면 회귀식의 계수를 재추정하여 산출 도구(\subsecref{subsec:practical-regression-tool})를
갱신해야 한다.}

\chg{요컨대 이중채널 감쇠 메커니즘은 ``일상에서는 관측성 투자, 극단에서는 불확실성 관리''라는 이중축 자원
배분 원칙과, Tail 시나리오를 별도로 설계하는 스트레스 테스트 기준을 동시에 제공한다. 이는 단일 지표나 평균
중심 의사결정으로는 포착할 수 없는, 분포 환경에 따른 입체적 위기 대응을 가능하게 하는 본 논문의 독자적
실무 기여이다.}


% ===================================================================
% 5.3.3 η 기반 평상시 성숙도·위기시 에스컬레이션 [구 성숙도 + 에스컬레이션 통합]
% ===================================================================
\subsection{\chg{$\eta$ 기반 평상시 성숙도와 위기시 에스컬레이션}}
\label{subsec:practical-bcm-std}        % 구 라벨 호환(성숙도)
\label{subsec:practical-decision}       % 구 라벨 호환(의사결정)
\label{subsec:escalation-thresholds}    % 구 라벨 호환
\label{subsec:eta-dual-use}             % 구 라벨 호환

\begin{p5add}
\chg{본 항은 단일 지표 $\eta = \DRTO/R_0$가 운영 모드에 따라 평상시 성숙도 진단 지표와 위기시 에스컬레이션
대응 지표의 두 역할을 수행함을 보이고, 그 평상시와 위기시 활용 및 선순환 구조를 제시한다. 평상시에는 $\eta$
변동이 주로 관측성 채널에 의해 결정되므로 BCMS 성숙도를 5단계 모형으로 진단하고, 위기시에는 $\eta$ 변동이
불확실성 채널에 의해 지배되므로 Green/Yellow/Orange/Red 4단계 에스컬레이션을 자동 트리거한다. 또한 데이터
수집과 모니터링 대시보드, 사후 분석으로 구성되는 BCM 운영 통합 도구가 함께 제공된다.}
\end{p5add}

\begingroup
  \renewcommand{\section}[1]{}
  \renewcommand{\subsection}[1]{}
  \renewcommand{\paragraph}[1]{\textbf{#1}\par}

  % η의 평상시·위기시 이중 활용 개념(overview) + 성숙도 5단계(1.1) + BCM 운영통합(1.2-4)
  
$\eta = \DRTO / R_0$는 단일 지표이지만, 조직의 운영 상황에 따라
두 가지 서로 다른 역할을 수행한다.
이러한 이중 활용이 가능한 근거는 $\eta$를 구성하는 두 채널\chg{(}관측성($E_{O,bnd}$)과
불확실성($E_{U,bnd}$)\chg{)}의 \textbf{지배적 기여 요인이 운영 모드에 따라 달라지기}
때문이다.
평상시에는 \chg{대규모 재난과 사고}가 부재하므로 불확실성이 \chg{일상적이고 가우시안}
수준에 머물러 $\eta$의 변동을 주로 관측성 역량이 결정한다.
반면 위기시에는 이미 \chg{사고와 재난}이 발생하여 규모를 알 수 없는 큰 불확실성이
현존하며, 이를 얼마나 빠르게 정량화하여 $\eta$에 반영하느냐가 대응의 핵심이 된다.
\p{[}표~\vref{tab:eta-dual-use}\p{]}는 이 \chg{평상시와 위기시} 활용 체계를 정리한 것이다.

\begin{table}[htbp]
\centering
\caption{$\eta$ 지표의 \chg{평상시·위기시} 활용 체계}
\label{tab:eta-dual-use}
\small
\begin{tabularx}{\textwidth}{l X X}
\toprule
\textbf{구분} & \textbf{평상시 (Normal Mode)} & \textbf{위기시 (Crisis Mode)} \\
\midrule
\textbf{역할} & BCMS 성숙도 평가 지표 & 에스컬레이션 대응 지표 \\
\textbf{핵심 질문} & ``조직의 복구 역량이 어느 수준인가?'' & ``현재 상황이 얼마나 심각한가?'' \\
\textbf{시간축} & 중장기 (월간/분기별 추이) & 실시간 (분$\sim$시간 단위) \\
\textbf{주요 활동} & 진단, 투자 계획, 벤치마크 & 등급 판정, 보고, 즉각 대응 \\
\textbf{대상자} & BCM 관리자, 경영진 & 복구 팀, 현장 담당자 \\
\textbf{산출물} & 성숙도 등급, 개선 로드맵 & 에스컬레이션 등급, 대응 조치 \\
\bottomrule
\end{tabularx}
\end{table}

\p{[}그림~\vref{fig:eta-dual-use-diagram}\p{]}는 $\eta$의 \chg{평상시와 위기시} 이중 활용과
BCMS 선순환 구조를 종합적으로 도식화한 것이다.
평상시 성숙도 평가(좌측)와 위기시 에스컬레이션 대응(우측)이
\chg{장애 발생과 복구 완료}를 통해 전환되며,
사후 분석과 역량 환류가 상단의 진단--향상--대응--환류 선순환을 형성한다.

\begin{figure}[htbp]
\centering
\resizebox{\textwidth}{!}{%
\begin{tikzpicture}[
  arr/.style={-{Stealth[length=2.5mm, width=2mm]}, line width=1.2pt},
  arrlbl/.style={font=\scriptsize\bfseries, fill=white, inner sep=1.5pt, rounded corners=1.5pt},
  sbox/.style={draw, rounded corners=2pt, semithick, minimum width=20mm,
    minimum height=9mm, align=center, font=\scriptsize},
]

%% ============================================================
%% 상단 박스: BCMS 선순환
%% ============================================================
\node[draw, rounded corners=6pt, line width=1pt,
  fill=green!6, draw=green!50!black,
  minimum width=120mm, minimum height=24mm] (topbox) at (0, 3.2) {};

\node[font=\small\bfseries, text=black] at (0, 4.15) {BCMS 선순환};

\node[sbox, fill=blue!12, draw=blue!50]  (s1) at (-3.7, 3.2) {\textbf{1. 진단}\\$\eta$ 추이};
\node[sbox, fill=blue!12, draw=blue!50]  (s2) at (-1.25, 3.2) {\textbf{2. 향상}\\$O\!\uparrow\; U\!\downarrow$};
\node[sbox, fill=red!12, draw=red!50]    (s3) at (1.25, 3.2)  {\textbf{3. 대응}\\에스컬레이션};
\node[sbox, fill=yellow!18, draw=yellow!60!black] (s4) at (3.7, 3.2) {\textbf{4. 환류}\\이력$\to$KPI};

\draw[-{Stealth[length=1.8mm]}, thin, blue!60] (s1) -- (s2);
\draw[-{Stealth[length=1.8mm]}, thin, red!50] (s2) -- (s3);
\draw[-{Stealth[length=1.8mm]}, thin, yellow!60!black] (s3) -- (s4);
\draw[-{Stealth[length=1.8mm]}, thin, green!50!black]
  (s4.south) -- ++(0,-0.28) -| (s1.south)
  node[pos=0.5, below, font=\scriptsize\bfseries, text=black,
    fill=green!6, inner sep=1.5pt, rounded corners=1.5pt] {$\eta$ 선순환};

%% ============================================================
%% 좌하 박스: 평상시 BCMS 성숙도
%% ============================================================
\node[draw, rounded corners=6pt, line width=1pt,
  fill=blue!6, draw=blue!60,
  minimum width=56mm, minimum height=38mm] (leftbox) at (-5.0, -1.0) {};

\node[font=\small\bfseries, text=black] at (-5.0, 0.55) {평상시: BCMS 성숙도 평가};

\node[anchor=center, inner sep=0pt] at (-5.0, -0.85) {%
  \scalebox{0.7}{%
  \begin{tabular}{c c c l}
  \toprule
  \textbf{단계} & \textbf{등급} & \textbf{$\eta$} & \textbf{수준} \\
  \midrule
  5 & 최적화 & $< 0.6$ & \textit{Optimized} \\
  4 & 관리 & $0.6$--$0.8$ & \textit{Managed} \\
  3 & 정의 & $0.8$--$1.0$ & \textit{Defined} \\
  2 & 반복 & $1.0$--$1.5$ & \textit{Repeatable} \\
  1 & 초기 & $\geq 1.5$ & \textit{Initial} \\
  \bottomrule
  \end{tabular}}%
};

\node[font=\scriptsize\bfseries, text=black] at (-5.0, -2.35) {%
  개선 두 축: 관측성$\uparrow$ / 불확실성$\downarrow$%
};

%% ============================================================
%% 우하 박스: 위기시 에스컬레이션
%% ============================================================
\node[draw, rounded corners=6pt, line width=1pt,
  fill=red!6, draw=red!60,
  minimum width=56mm, minimum height=38mm] (rightbox) at (5.0, -1.0) {};

\node[font=\small\bfseries, text=black] at (5.0, 0.55) {위기시: 에스컬레이션 대응};

\node[anchor=center, inner sep=0pt] at (5.0, -0.85) {%
  \scalebox{0.7}{%
  \begin{tabular}{c c l l}
  \toprule
  \textbf{$\eta$} & \textbf{등급} & \textbf{보고} & \textbf{대응} \\
  \midrule
  $< 0.8$ & Green & 팀 자체 & 정상 복구 \\
  $0.8$--$1.0$ & Yellow & 팀장 & 모니터링 강화 \\
  $1.0$--$1.5$ & Orange & 부서장 & 자원 추가 투입 \\
  $\geq 1.5$ & Red & 경영진 & BCP 전면 가동 \\
  \bottomrule
  \end{tabular}}%
};

\node[font=\scriptsize\bfseries, text=black] at (5.0, -2.35) {%
  국가 위기경보 연동: 관심--주의--경계--심각%
};

%% ============================================================
%% 3개 박스 연결 화살표
%% ============================================================
\draw[arr, red!50]
  ([yshift=-4mm]leftbox.east) -- ([yshift=-4mm]rightbox.west)
  node[arrlbl, midway, below=1.5pt, text=black] {장애$\cdot$재난 발생 $\to$ 위기 모드 전환};
\draw[arr, blue!50]
  ([yshift=4mm]rightbox.west) -- ([yshift=4mm]leftbox.east)
  node[arrlbl, midway, above=1.5pt, text=black] {복구 완료 $\to$ 평상 모드 복귀};
\draw[arr, green!50!black]
  ([xshift=6mm]rightbox.north) -- ([xshift=17mm]topbox.south)
  node[arrlbl, midway, right=1.5pt, text=black] {사후 분석};
\draw[arr, blue!60]
  ([xshift=-17mm]topbox.south) -- ([xshift=-6mm]leftbox.north)
  node[arrlbl, midway, left=1.5pt, text=black] {역량 환류};

%% ============================================================
%% 하단 강조 문구
%% ============================================================
\node[font=\small\bfseries, text=black, align=center,
  draw=red!50, rounded corners=3pt, fill=red!6, inner sep=4pt,
  line width=0.6pt] at (0, -3.5) {%
  $\eta$ 하나로 평상시 ``얼마나 준비되었는가'' + 위기시 ``얼마나 심각한가'' 통합 측정
  $\;\to\;$ \textcolor{black}{\bfseries 데이터 기반 BCMS 선순환}%
};

\end{tikzpicture}%
}% end resizebox
\caption{$\eta$의 \chg{평상시와 위기시} 이중 활용 체계와 BCMS 선순환 구조}
\label{fig:eta-dual-use-diagram}
\end{figure}

  %% --- 평상시: BCMS 성숙도 평가 ---
\subsubsection{\tmp{평상시 BCMS 성숙도 평가}}
\label{subsubsec:maturity-model}

평상시에는 \chg{대규모 사고와 재난}이 발생하지 않은 상태이므로,
불확실성은 일상적 업무 변동(가우시안 수준)에 한정되어 $E_{U,bnd}$의 기여가 작다.
따라서 $\eta$의 변동은 주로 \textbf{관측성 역량}($E_{O,bnd}$)에 의해 결정되며,
$\eta$의 중장기 추이를 추적하면 조직이 모니터링 체계를 얼마나 잘 구축하고
있는지\chg{,} 즉 복구 역량의 성숙도를 정량적으로 진단할 수 있다.
$\eta$가 낮을수록 관측성 역량이 우수한 성숙한 조직으로 해석되며,
이에 기반하여 BCMS 성숙도를 5단계로 구분하는 모형을 \p{[}표~\vref{tab:maturity-5level}\p{]}에서 제안한다.

각 임계값은 본 연구의 시뮬레이션 결과에 기반한다\chg{.}
$\eta = 0.6$은 C1 조건의 하위 \erf{25}{5}\% 수준으로 관측성이 극대화된 상태,
$\eta = 0.8$은 C1 평균 $\bar{\eta} = 0.853$ 이하로 $\DRTO < R_0$가 안정적으로 유지되는 상태,
$\eta = 1.0$은 $\DRTO = R_0$ 기준선으로 관측성과 불확실성이 상쇄되는 지점,
$\eta = 1.5$는 불확실성이 지배적이어서 $\DRTO$가 $R_0$의 1.5배를 초과하는 상태이다.

\begin{table}[htbp]
\centering
\caption{$\eta$ 기반 BCMS 성숙도 5단계 모형}
\label{tab:maturity-5level}
\small
\begin{tabularx}{\textwidth}{c c c l X}
\toprule
\textbf{단계} & \textbf{등급} & \textbf{$\eta$ 범위} & \textbf{수준} & \textbf{특징} \\
\midrule
5 & 최적화 & $\eta < 0.6$ & \textit{Optimized} &
  관측성 극대화, 불확실성 최소화. 자동화된 모니터링·복구 \\
\addlinespace
4 & 관리 & $0.6 \leq \eta < 0.8$ & \textit{Managed} &
  정량적 KPI 기반 BCM 운영. $\DRTO < R_0$ 안정 유지 \\
\addlinespace
3 & 정의 & $0.8 \leq \eta < 1.0$ & \textit{Defined} &
  공식 BCM 프로세스 수립. 관측성과 불확실성 \m{상쇄($\DRTO \approx R_0$)} \\
\addlinespace
2 & 반복 & $1.0 \leq \eta < 1.5$ & \textit{Repeatable} &
  기본 절차 존재하나 불확실성 우세. $\DRTO > R_0$ 빈번 \\
\addlinespace
1 & 초기 & $\eta \geq 1.5$ & \textit{Initial} &
  BCM 체계 미비. 극단적 불확실성 노출 \\
\bottomrule
\end{tabularx}
\end{table}

성숙도를 향상(즉, $\eta$를 감소)시키는 경로는 두 축으로 구분된다.
\chg{첫째, 관측성 강화($E_{O,bnd}$ 증가로 $\eta$ 감소)는 실시간 모니터링 체계 구축, APM/SIEM 연동, 장애
원인 분석(RCA) 체계화를 포함하며, 본 연구의 C1 조건에서 $k_O$ 계수가 관측성의 실질적 효과를 결정하므로
$k_O$가 높은 영역에 우선 투자하는 것이 효율적이다. 둘째, 불확실성 저감($E_{U,bnd}$ 감소로 $\eta$ 감소)은
복구 절차서 최신 유지, 정기 BCM 훈련, 공급망 대체 방안 확보를 포함하며, 특히 Tail 영역(P85.8 초과)에서는
파레토 불확실성이 지배적이므로 블랙스완 대비 계획이 성숙도 향상의 핵심이 된다.}

각 성숙도 단계별 구체적 전환 전략, 위기 시 정책 활용 방안, 도입 로드맵은
아래에 상세히 기술하며, 현업 담당자용 자가 진단 체크리스트(15개 항목)는
\appchref{app:eta-dual-use-detail}에 수록한다.

%% [분량보존 복원] η 평상시·비상시 활용 상세 (구 부록 H.2~H.4): 향상전략·위기정책·로드맵
\subsubsection{성숙도 단계별 향상 전략}
\label{app:maturity-strategy}

각 성숙도 단계에서 다음 단계로 향상하기 위한 구체적 전략을 기술한다.
핵심은 관측성 강화($E_{O,bnd}$ 증가)와 불확실성 저감($E_{U,bnd}$ 감소)의
두 레버를 단계에 맞게 적용하는 것이다.

\chg{1단계(Initial)에서 2단계(Repeatable)로의 기반 구축 단계에서는 기본 모니터링 도구 도입(가용성
모니터링, 기본 알림), BCM 정책 문서 및 복구 절차서(runbook) 초안 작성, BIA 최초 수행을 통한 RTO
기준값($R_0$) 설정, RPO 기준값 설정, BCM 전담 책임자 지정을 수행하며, 목표는 $\eta < 1.5$ 안정화이다.}

\chg{2단계(Repeatable)에서 3단계(Defined)로의 체계화 단계에서는 관측성 3축(로그, 메트릭, 트레이스) 중
2개 이상 확보, 장애 대응 절차 표준화 및 정기 훈련(연 2회 이상), MTTR 측정과 추적 시작을 통한
$O_{\text{raw}}$ 데이터 수집, 에스컬레이션 기준 및 보고 체계 문서화를 수행하며, 목표는 $\eta < 1.0$($\DRTO \leq R_0$)이다.}

\chg{3단계(Defined)에서 4단계(Managed)로의 정량화 단계에서는 실시간 $\DRTO/R_0$ 모니터링을 위한 $\eta$
대시보드 구축, $k_O$ 변화 추적을 통한 관측성 투자 ROI 측정, 불확실성 원인별 분류 체계 수립(가우시안 vs
파레토 유형), 사후 분석을 통한 절차 개선 피드백 루프 정착을 수행하며, 목표는 $\eta < 0.8$ 유지(에스컬레이션
Green 영역)이다.}

\chg{4단계(Managed)에서 5단계(Optimized)로의 최적화 단계에서는 AIOps 기반 자동 장애 탐지와 진단, 복구
체계 구축, 카오스 엔지니어링(Chaos Engineering) 정기 수행, 향후 $\eta$ 추이 예측 및 선제 대응을 위한
$\eta$ 예측 모형 운영, 극단 시나리오(C3) 전용 대응을 위한 CVaR$_{99}$ 기반 자원 배분을 수행하며, 목표는
$\eta < 0.6$ 달성 및 유지이다.}


\subsubsection{위기시 정책 활용 방안}
\label{app:crisis-policy}

\chg{국가 위기경보 연동 $\eta$ 임계값 자동 조정은}
국가 위기경보가 상향될 때 조직 내부의 에스컬레이션 민감도를
자동으로 높여 선제 대응하는 정책이다.
\nref{위기경보 4단계에 따른 $\eta$ 임계값 조정 규칙은 \tabref{tab:app-policy-adjustment}와 같다.}

\begin{table}[htbp]
\centering
\caption{국가 위기경보 연동 $\eta$ 임계값 조정 정책}
\label{tab:app-policy-adjustment}
\small
\newcolumntype{W}{>{\raggedright\arraybackslash}p{0.42\textwidth}}
\begin{tabularx}{\textwidth}{c c X W}
\toprule
\textbf{위기경보} & \textbf{임계값 조정} & \textbf{근거} & \textbf{조정 후 기준} \\
\midrule
관심(Blue) & 기본값 유지 &
  일상 운영 &
  Green $\eta < 0.8$,\; Yellow $0.8 \leq \eta < 1.0$,\; Orange $1.0 \leq \eta < 1.5$,\; Red $\eta \geq 1.5$ \\
\addlinespace
주의(Yellow) & 각 $-0.1$ &
  외부 위험 증가 &
  Green $\eta < 0.7$,\; Yellow $0.7 \leq \eta < 0.9$,\; Orange $0.9 \leq \eta < 1.4$,\; Red $\eta \geq 1.4$ \\
\addlinespace
경계(Orange) & 각 $-0.2$ &
  위기 가능성 높음 &
  Green $\eta < 0.6$,\; Yellow $0.6 \leq \eta < 0.8$,\; Orange $0.8 \leq \eta < 1.3$,\; Red $\eta \geq 1.3$ \\
\addlinespace
심각(Red) & 즉시 Red &
  대규모 위기 &
  $\eta$와 무관하게 BCP 전면 가동, 전 등급 Red 적용 \\
\bottomrule
\end{tabularx}
\end{table}


\chg{위기 유형별 $\eta$ 가중 모형은}
재난 유형에 따라 $\eta$ 구성 요소(관측성 vs 불확실성)에 가중치를 부여한다.
\nref{재난 4유형별 $O\!:\!U$ 가중 비율과 대응 전략 초점은 \tabref{tab:app-crisis-type-weight}에 정리하였다.}

\begin{table}[htbp]
\centering
\caption{위기 유형별 대응 전략 가중 모델}
\label{tab:app-crisis-type-weight}
\small
\begin{tabularx}{\textwidth}{l c c X}
\toprule
\textbf{위기 유형} & \textbf{$O$ 비중} & \textbf{$U$ 비중} & \textbf{대응 전략 초점} \\
\midrule
사이버 공격 & 높음 (70\%) & 보통 (30\%) &
  SIEM/EDR 강화, 탐지·진단 역량 집중 \\
\addlinespace
자연재해 & 보통 (40\%) & 높음 (60\%) &
  물리적 피해 불확실성 대비, 대체 사이트 확보 \\
\addlinespace
감염병 & 보통 (50\%) & 높음 (50\%) &
  원격 운영 역량 + 인력 가용성 불확실성 동시 대응 \\
\addlinespace
공급망 중단 & 낮음 (30\%) & 높음 (70\%) &
  외부 불확실성 지배. 다중 공급원 확보, 재고 완충 \\
\bottomrule
\end{tabularx}
\end{table}


\chg{에스컬레이션 이력 기반 성숙도 KPI로서,}
위기시 에스컬레이션 이력은 평상시 성숙도 개선의 데이터로 환류된다.
\chg{주요 KPI는 네 가지이다. Green 체류 비율은 전체 운영 시간 대비 $\eta < 0.8$ 비율로 성숙도가 높을수록
증가하며, 에스컬레이션 빈도는 월간 Yellow 이상 발생 횟수로 감소 추세가 성숙도 향상을 의미하고, 복귀
시간은 Orange/Red에서 Green으로의 복귀 소요 평균 시간으로 단축이 복구 역량 향상을 뜻하며, Red 도달률은
위기(Red) 등급 발생 빈도로 감소가 극단 대비 역량 향상을 나타낸다.}

이러한 KPI를 월간/분기별로 경영진에 보고함으로써
BCM 투자의 효과를 정량적으로 증명할 수 있다.


\subsubsection{$\eta$ \chg{평상시와 비상시} 활용 체계 도입 로드맵}
\label{app:eta-roadmap}

\nref{$\eta$ 기반 BCM 체계의 단계별 도입 일정과 주요 활동은 \tabref{tab:app-roadmap}에 정리하였다.}

\begin{table}[htbp]
\centering
\caption{$\eta$ 평상시·비상시 활용 체계 도입 로드맵}
\label{tab:app-roadmap}
\small
\begin{tabularx}{\textwidth}{c l X}
\toprule
\textbf{단계} & \textbf{기간} & \textbf{주요 활동} \\
\midrule
준비 & 1--2개월 &
  BIA 수행 $\to$ $R_0$ 설정. 관측성 데이터 수집 시작.
  $\eta$ 산출 파일럿 구축. 자가 진단 최초 실시 \\
\addlinespace
파일럿 & 3--4개월 &
  핵심 시스템 1--2개에 $\eta$ 대시보드 적용.
  에스컬레이션 기준 시범 운영. 국가 위기경보 연동 테스트 \\
\addlinespace
확산 & 5--8개월 &
  전사 확대 적용. 성숙도 자가 진단 분기별 정기화.
  에스컬레이션 이력 KPI 보고 시작 \\
\addlinespace
최적화 & 9--12개월 &
  AIOps 연동. $\eta$ 예측 모형 운영.
  위기 유형별 가중 모델 적용. 산업별 벤치마크 비교 \\
\bottomrule
\end{tabularx}
\end{table}






  %% -------------------------------------------------------------------
\subsection{\vrev{BCM 운영 통합 프레임워크}}
\label{subsec:bcm-integration}
%% -------------------------------------------------------------------

2단계 프레임워크와 에스컬레이션 체계를 기존 BCMS에 통합하기 위한
구체적 운영 절차를 다음과 같이 제시한다.

\subsubsection{\tmp{데이터 수집 체계}}

\inew{$\DRTO$ 모델의 두 입력 변수인 \textbf{관측성 $O(t)$} 와 \textbf{불확실성 $U(t)$} 의
수집은 각각 기존 \chg{운영과 보안} 모니터링 인프라와 \chg{장애 이력 및 위협 인텔리전스} 데이터의
연동을 통해 구현한다.}

\inew{\chg{관측성 지표 $O(t)$의 수집은 기존 모니터링 인프라와의 연동을 통해 다음과 같이 구현한다. 첫째, APM(Application Performance Monitoring) 연동으로 서비스 응답 시간, 오류율, 처리량을 종합하여 $O_{\mathrm{APM}} \in [0,1]$로 정규화하며, 가중 평균은 $O_{\mathrm{APM}} = w_1 \cdot R_{\mathrm{avail}} + w_2 \cdot (1-R_{\mathrm{error}}) + w_3 \cdot R_{\mathrm{perf}}$($\sum w_i = 1$)로 산출한다. 둘째, SIEM(Security Information and Event Management) 연동으로 보안 이벤트 심각도와 대응 상태를 기반으로 $O_{\mathrm{SIEM}} \in [0,1]$을 산출하되, 미대응 고위험 이벤트가 존재하면 감점한다. 셋째, 종합 관측성은 $O(t) = \alpha \cdot O_{\mathrm{APM}} + (1-\alpha) \cdot O_{\mathrm{SIEM}}$로 정의하며, $\alpha$는 조직 특성에 따라 예를 들어 $[0.5, 0.8]$에서 설정한다.}}

\inew{\chg{불확실성 지표 $U(t)$의 수집은 조직 내부의 과거 장애 이력과 외부 위협 인텔리전스 데이터의 연동을 통해 다음과 같이 구현한다. 첫째, 장애 이력 DB 연동으로 조직 자체의 과거 장애와 복구 이력(MTTR, RTO 초과 사례, 복구 소요 시간 분포)을 적합 분포(가우시안 또는 파레토)로 모델링하여 $U_{\mathrm{hist}} \in [0,\; U_{\max}]$을 산출하며, 분포 형태 판정에는 Hill 추정량~\citep{hill1975simple} 또는 Anderson-Darling 검정~\citep{anderson1954test}을 적용하여 C2(경꼬리) 또는 C3(중꼬리) 환경으로 자동 분기한다. 둘째, 위협 인텔리전스(CVE/MITRE ATT\&CK) 연동으로 자산 대장과 매핑된 CVE 공개 DB, EPSS 점수, MITRE ATT\&CK 위협 액터 활동 정보를 기반으로 미래 위험 강도 $U_{\mathrm{threat}} \in [0,\; U_{\max}]$을 산출하되, 미패치 고위험 CVE 수와 EPSS 가중 평균을 결합한다. 셋째, 종합 불확실성은 $U(t) = \beta \cdot U_{\mathrm{hist}} + (1-\beta) \cdot U_{\mathrm{threat}}$로 정의하며, $\beta$는 조직의 과거 데이터 풍부도와 위협 환경 변동성에 따라 $[0.5, 0.8]$에서 설정한다. 장애 이력이 풍부한 성숙 조직은 $\beta$가 높고, 신생 조직과 신규 위협 환경은 $\beta$가 낮다.}}

\inew{원시 변수 $O_{\mathrm{raw}}$, $U_{\mathrm{raw}}$, 가중치 $k_O$의 제1계층
하위 지표 조작화 예시와 가상 기업 라이프사이클 적용 시나리오는
\secref{app:operationalization}(관측성 $O_{\mathrm{raw}}$ 조작화 예시,
불확실성 $U_{\mathrm{raw}}$ 조작화 예시)에 참고용으로 수록되어 있다.}

\subsubsection{\tmp{$\eta$ 모니터링 대시보드}}

$\eta$ 비율의 실시간 시각화를 통해 BCMS 효과성을 지속적으로 측정한다.
대시보드의 핵심 구성요소는 \chg{실시간 $\eta$ 추이 차트(시계열),
에스컬레이션 등급 표시(Green/Yellow/Orange/Red),
$O(t)$와 $U(t)$의 개별 추이, CVaR$_{99}$ 현재값}이며,
이는 \nrev{ISO 22301:2019} 조항~9.1(모니터링, 측정, 분석, 평가)의
정량적 요구사항을 직접 충족한다.

\subsubsection{\tmp{사후 분석 프로세스}}

복구 완료 후 24--72시간 이내에 다음의 사후 분석을 수행한다.
\chg{먼저 실제 복구 시간 $T_{\mathrm{actual}}$과 $\DRTO$ 예측값을 비교하고,
절대 오차 $E = T_{\mathrm{actual}} - \DRTO_{\mathrm{final}}$과
상대 오차 $E_r = E / T_{\mathrm{actual}}$를 산출하며,
과대 및 과소 추정 원인을 분석하여 파라미터 교정을 권고하고,
장애 이력 DB 갱신 및 $U(t)$ 분포 재추정을 수행한다.}
이 프로세스는 \nrev{ISO 22301:2019} 조항~10.1(지속적 개선)과 직접 대응된다.




  % 위기시 에스컬레이션(3.1) + 선순환(3.2)
  %% --- 위기시: 에스컬레이션 대응 체계 ---
\subsubsection{\syncr{위기 시 에스컬레이션 대응 체계}}
\label{subsubsec:escalation-detail}

위기시에는 \chg{사고와 재난}이 이미 발생한 상태이므로,
규모를 알 수 없는 큰 불확실성이 현존한다.
이때 $\eta$는 이 불확실성의 크기를 실시간으로 반영하는 위험 수준 지표가 되며,
조직이 불확실성을 얼마나 빠르게 정량화하여 $\eta$에 반영하느냐가
적시 대응의 관건이 된다.
$\eta$ 값에 따라 4단계 에스컬레이션 체계가 자동 결정된다.
\p{[}표~\vref{tab:ch6-eta-escalation}\p{]}는 $\eta$ 구간별 대응 체계를 정의하며,
\p{[}그림~\vref{fig:ch6-escalation-zones}\p{]}는 이를 시각화한 것이다.

\begin{table}[htbp]
\centering
\caption{$\eta$ 기반 에스컬레이션 임계값 및 대응 체계}
\label{tab:ch6-eta-escalation}
\small
\begin{tabularx}{\textwidth}{c c l l X}
\toprule
\textbf{단계} & \textbf{$\eta$ 범위} & \textbf{등급} & \textbf{보고 대상} & \textbf{대응 조치} \\
\midrule
1 & $\eta < 0.8$ & \m{안정(Green)} & 복구 팀 자체 & 정상 복구 절차 수행; 관측성에 의한 $\DRTO$ 단축 효과 활성 \\
\addlinespace
2 & $0.8 \leq \eta < 1.0$ & \m{주의(Yellow)} & 팀장 보고 & 불확실성 증가 모니터링; $\DRTO$ 갱신 주기 \m{단축(15분)}; 추가 자원 대기 \\
\addlinespace
3 & $1.0 \leq \eta < 1.5$ & \m{경고(Orange)} & 부서장 보고 & $\DRTO > R_0$로 기준 RTO 초과; 추가 자원 투입; 외부 지원 요청; 이해관계자 긴급 통보 \\
\addlinespace
4 & $\eta \geq 1.5$ & \m{위기(Red)} & 경영진 보고 & 위기관리 모드 전환; 사업 연속성 비상 계획(BCP) 전면 가동; CVaR$_{99}$ 기준 최악 시나리오 대비 \\
\bottomrule
\end{tabularx}
\end{table}

\chg{에스컬레이션 트리거는 $\eta$ 급등(30\% 이상 증가), 조건 상향 전환(C2 $\to$ C3), 에스컬레이션 임계값
접근($\eta > 1.4$), 불확실성 급등($U > 0.8$)의 네 가지이다.}

\begin{figure}[htbp]
\centering
\resizebox{0.95\textwidth}{!}{%
\begin{tikzpicture}
% 배경 영역
\fill[green!15]  (0,0)   rectangle (3.2,3.5);
\fill[yellow!20] (3.2,0) rectangle (6.4,3.5);
\fill[orange!20] (6.4,0) rectangle (11.2,3.5);
\fill[red!20]    (11.2,0) rectangle (14.4,3.5);

% 경계선
\draw[thick, black] (3.2,0) -- (3.2,3.5);
\draw[thick, black] (6.4,0) -- (6.4,3.5);
\draw[thick, black] (11.2,0) -- (11.2,3.5);

% 축
\draw[thick, -{Stealth[length=3mm]}] (0,0) -- (15,0) node[right, font=\small] {$\eta$};
\draw[thick] (0,0) -- (0,3.5);

% η 값 레이블
\node[below, font=\footnotesize] at (0, -0.15) {$0$};
\node[below, font=\footnotesize\bfseries] at (3.2, -0.15) {$0.8$};
\node[below, font=\footnotesize\bfseries] at (6.4, -0.15) {$1.0$};
\node[below, font=\footnotesize\bfseries] at (11.2, -0.15) {$1.5$};

% 영역 레이블 (상단)
\node[font=\small\bfseries, green!50!black] at (1.6, 3.0) {안정};
\node[font=\small\bfseries, yellow!60!black] at (4.8, 3.0) {주의};
\node[font=\small\bfseries, orange!60!black] at (8.8, 3.0) {경고};
\node[font=\small\bfseries, red!60!black] at (12.8, 3.0) {위기};

% 영문 등급
\node[font=\scriptsize, green!40!black] at (1.6, 2.5) {Green};
\node[font=\scriptsize, yellow!50!black] at (4.8, 2.5) {Yellow};
\node[font=\scriptsize, orange!50!black] at (8.8, 2.5) {Orange};
\node[font=\scriptsize, red!50!black] at (12.8, 2.5) {Red};

% 핵심 조건 설명
\node[font=\scriptsize, align=center] at (1.6, 1.5)
  {관측성에 의한\\$\DRTO$ 단축\\$E_{O,bnd}$ 지배};
\node[font=\scriptsize, align=center] at (4.8, 1.5)
  {불확실성 증가\\관측성 효과 상쇄\\$E_{O,bnd} + E_{U,bnd} \approx 0$};
\node[font=\scriptsize, align=center] at (8.8, 1.5)
  {$\DRTO > R_0$\\불확실성 프리미엄\\C2 기반 자원 추가};
\node[font=\scriptsize, align=center] at (12.8, 1.5)
  {파레토 활성화\\$0.52$배 감쇠\\BCP 전면 가동};

% R₀ 기준선 표시
\draw[dashed, blue!60, thick] (6.4, -0.5) -- (6.4, -0.9)
    node[below, font=\scriptsize, blue!60] {$\eta = 1.0$: $\DRTO = R_0$ 기준};

\end{tikzpicture}
}% end resizebox
\caption{$\eta$ 기반 에스컬레이션 임계값 \m{다이어그램(Green / Yellow / Orange / Red)}}
\label{fig:ch6-escalation-zones}
\end{figure}

에스컬레이션 임계값의 설정 근거는 다음과 같다.
$\eta = 0.8$은 C1 조건에서의 평균 $\eta = 0.853$보다 낮은 값으로,
관측성의 단축 효과가 충분히 발현되어 추가 대응이 불필요한 안정 상태를 의미한다.
$\eta = 1.0$은 $\DRTO = R_0$, 즉 DRTO가 기준 RTO와 동일한 지점으로,
이 이상에서는 불확실성 프리미엄에 의해 복구 시간이 기준을 초과한다.
$\eta = 1.5$는 $\DRTO = 1.5 \times R_0 = 9.0$\,h에 대응하며,
이 이상에서는 시그모이드 바운딩의 포화 영역에 진입하여
위기관리 모드 전환이 필요하다.
이러한 단계별 에스컬레이션은 SRE의 에러 예산 소진율 기반
인시던트 대응 체계~\citep{beyer2016sre}와 유사한 구조를 가지나,
$\eta$가 관측성과 불확실성의 결합 효과를 단일 지표로 포착한다는 점에서
BCM 맥락에 최적화된 확장이다.

\chg{나아가 2단계 프레임워크와의 복합 에스컬레이션이 가능하다.}
$\eta$ 기반 4단계 등급은 \subsubsecref{subsec:two-stage-framework}의
2단계 프레임워크(Core C2 + Tail C3)와 결합하여 더욱 정교한 대응이 가능하다.
동일한 $\eta$ 등급이라도 불확실성 분포가 가우시안(C2)인지
파레토(C3)인지에 따라 위험의 성격이 질적으로 다르기 때문이다.
예를 들어, $\eta = 1.3$(Orange 등급)이 Core 영역(C2)에서 관측되면
가우시안 불확실성의 일시적 증가로 해석되어 추가 자원 투입과 모니터링 강화가
1차 대응이 된다. 그러나 동일 $\eta$가 Tail 영역(C3)에서 관측되면
파레토 꼬리 위험이 활성화된 상태이므로, $k_O$ 계수의 $0.52$배 감쇠로 인해
관측성 효과가 제한적이며, CVaR$_{99}$가 $+24.0\%$ 높은 극단 시나리오에
대비한 선제적 BCP 점검이 필요하다.

이를 운영 규칙으로 정리하면 다음과 같다.
\chg{먼저 Core 영역(P85.8 이하)에서는 $\eta$ 등급에 따른 표준 에스컬레이션 절차를 적용하고,
Tail 영역(P85.8 초과)에서는 동일 $\eta$ 등급이라도 한 단계 높은 대응 수준을
적용하며(예: Tail에서의 Orange $\to$ Red 수준 대응),
Core에서 Tail로의 조건 전환(C2 $\to$ C3) 자체를
에스컬레이션 트리거로 취급한다.}
이러한 복합 에스컬레이션은 $\eta$라는 단일 축의 한계를 보완하여,
불확실성의 분포 특성까지 반영한 입체적 위기 대응을 가능하게 한다.

\chg{또한 국가 위기경보 체계와도 대응된다.}
본 연구의 $\eta$ 기반 에스컬레이션 4단계는 ``재난 및 안전관리 기본법''에 근거한
\textbf{국가 위기경보 4단계}(관심--주의--경계--심각)와 구조적으로 대응된다.
\p{[}표~\vref{tab:eta-crisis-alert-mapping}\p{]}은 양 체계의 매핑을 보여준다.

\begin{table}[htbp]
\centering
\caption{$\eta$ 기반 에스컬레이션과 국가 위기경보 4단계 대응}
\label{tab:eta-crisis-alert-mapping}
\small
\begin{tabularx}{\textwidth}{c c c c c X}
\toprule
\textbf{$\eta$ 범위} & \textbf{DRTO 등급} & \textbf{색상} & \textbf{위기경보} & \textbf{색상} & \textbf{대응 원칙의 공통점} \\
\midrule
$\eta < 0.8$ & 안정 & Green & 관심 & 파랑 &
  징후 감지 수준; 정상 운영 유지, 모니터링 강화 \\
\addlinespace
$0.8 \leq \eta < 1.0$ & 주의 & Yellow & 주의 & 노랑 &
  위험 가능성 증가; 자원 점검, 갱신 주기 단축 \\
\addlinespace
$1.0 \leq \eta < 1.5$ & 경고 & Orange & 경계 & 주황 &
  기준 초과 확인; 추가 자원 투입, 이해관계자 통보 \\
\addlinespace
$\eta \geq 1.5$ & 위기 & Red & 심각 & 빨강 &
  총력 대응; BCP 전면 가동, 경영진 의사결정 \\
\bottomrule
\end{tabularx}
\end{table}

두 체계는 \chg{위험 강도에 따른 단계적 격상(graduated response),
색상코드 기반 직관적 커뮤니케이션,
각 단계별 보고 대상 및 대응 수준의 명확한 구분}이라는
공통 설계 원칙을 공유한다.
차이점은 국가 위기경보가 질적 판단에 기반하는 반면,
$\eta$ 기반 에스컬레이션은 $\DRTO / R_0$라는 \textbf{정량적 지표}에
의해 자동으로 등급이 결정된다는 점이다.
이러한 정량적 트리거는 주관적 판단 편향을 배제하고
일관된 의사결정을 지원하며, 기존 위기경보 체계에 보완적으로
통합될 수 있다.


  
%% --- 선순환 구조 ---
\subsubsection{\syncr{진단에서 향상, 대응, 환류로 이어지는 선순환 구조}}
\label{subsubsec:virtuous-cycle}

$\eta$의 \chg{평상시와 위기시} 활용이 만드는 핵심 가치는 \textbf{선순환}이다.
평상시의 성숙도 진단과 역량 향상이 위기시의 대응력을 높이고,
위기시의 에스컬레이션 이력\chg{(}Green 체류 비율,
복귀 시간, Red 도달률 등\chg{)}이 평상시의 성숙도 개선 데이터로 환류된다.
이 순환의 공통 언어가 $\eta$이며,
BCMS가 ``문서 기반 형식적 체계''에서 ``데이터 기반 실행적 체계''로
전환되는 것이 $\DRTO$ 프레임워크가 실무에 제공하는 핵심 가치이다.


\endgroup

\subsubsection{\syncr{실무 적용 가능성 예시}}
\label{subsubsec:practical-application-example}

\syncr{지금까지 제시한 $\eta$ 기반 \chg{평상시와 위기시} 이중 활용 체계가 실제 조직에
적용될 경우, 관측성($O_{\text{raw}}, k_O$)과 불확실성($U$)은 APM/SIEM 로그,
장애 이력 DB, CVE 공개 DB, 자산 대장 등 현장의 기존 운영 시스템에서
직접 관측되는 원시 지표로부터 산출 가능할 것이다.
참고를 위해 아래 \subsecref{app:operationalization}에 제1계층 하위 지표의
조작화 예시와 가상 기업의 4단계 라이프사이클
(진단→향상→대응→환류) 적용 가능 시나리오를 제시한다.
본 예시는 실무 적용 방향성을 보여주기 위한 참고용이며, 하위 지표의
구체적 \chg{구성과 가중치, 측정 도구}의 실증적 검증은 본 연구의 범위를 벗어나는
후속 연구 과제로 남겨둔다(\secref{sec:limitations} 참조).}



%% [본문 이관] 구 부록 J: 하위 변수 조작화 및 가상 기업 적용 예시 (참고용)
\subsection{\syncr{하위 변수 조작화 및 가상 기업 적용 예시}}
\label{app:operationalization}

\syncr{\subsubsection{하위 변수 조작화의 성격과 한계}}

\syncr{\textbf{하위 변수 조작화는 $\DRTO$ 프레임워크의 실무 적용 방향성을 보여주기 위한
참고용 예시이다.} 제시된 하위 지표의 구성, 가중치, 측정 도구 매핑,
가상 기업 수치는 \textbf{실증적으로 검증되지 않았으며}, 중견 제조업 일반
특성에서 추정한 \textbf{조건적 시나리오}(illustrative case~\citep{yin2018case})에
해당한다. 이를 규범적 가이드라인 또는 검증된 방법론으로 해석하지 말 것을 당부한다.}

\syncr{본문에서 제시한 $\DRTO$ \chg{모델과 시뮬레이션 결과, 전문가 검증}은 본 연구의
공식적 기여이나, 하위 변수 조작화와 시나리오는 \textbf{향후 연구를 위한
탐색적 제안}에 그친다. 각 하위 지표의 심리측정학적 타당도, 가중치의
산업별 튜닝, 원시 지표의 측정 도구별 \chg{정확도와 편향}, 현장 데이터 기반
\chg{수렴 및 예측} 타당도는 본 연구의 범위를 벗어나며 후속 연구에서 독립적으로
검증되어야 한다.}

\syncr{\subsubsection{조작화의 필요성과 정지점 원칙}}

\syncr{$\DRTO$ 모델의 주 변수인 관측성($O_{\text{raw}}$, $k_O$)과 불확실성($U$)은
수식적으로는 정의되어 있으나, 실제 조직에서 이들을 어떻게 산출할 것인가는
추가 조작적 정의가 필요하다. 조작화는 원칙적으로 하위 변수의 하위 변수로
무한 분해될 수 있으므로, 어느 수준에서 정지할 것인가의 \textbf{정지점
(stopping point)} 문제가 발생한다.}

\syncr{하위 변수 조작화는 측정이론의 고전적 원칙을 따라, 해당 하위 지표가 현장의
\textbf{기존 운영 시스템}(APM 도구·SIEM 도구·장애 이력 데이터베이스·CVE 공개 데이터베이스·자산 대장 등)\textbf{에서 직접 \chg{관측되고 기록되는} 원시값}에
도달하면 그 이하의 분해는 본 연구의 범위 밖으로 간주한다. 이는
\synct{bagozzi1998general}의 ``observable primitive'' 정지 규칙,
\synct{bollen1989structural}의 effect indicator 이론, 그리고
\nrev{ISO/IEC 25010:2011}의 원자 지표 표준화 원칙에 부합한다~\citep{iso25010}.}

\syncr{\chg{다만} 본 정지 규칙은 일반적 원칙을 따른 것이며,
개별 원시 지표의 심리측정학적 타당도는 측정 도구별로 별도 검증이 필요하다.}

\syncr{\subsubsection{관측성 $O_{\text{raw}}$ 조작화 예시}}

\syncr{관측성 원시값 $O_{\text{raw}} \in [0,1]$는 응용 성능 모니터링(APM)과 보안 \chg{정보 및 이벤트}
관리(SIEM)의 두 차원으로 분해 가능하다.}

\syncr{\chg{APM 차원은 $O_{\text{APM}} = w_1 R_{\text{avail}} + w_2 (1-R_{\text{error}}) + w_3 R_{\text{perf}}$로 정의한다.}}

\syncr{\chg{SIEM 차원은 $O_{\text{SIEM}} = v_1 (1-E_{\text{high}}) + v_2 (1-T_{\text{mttr}})$로 정의한다.}}

\syncr{\chg{두 차원의 합성은 $O_{\text{raw}} = \alpha \cdot O_{\text{APM}} + (1-\alpha) \cdot O_{\text{SIEM}}$이며, 본 예시에서 $\alpha = 0.7$로 가정한다.}}

\syncr{제1계층 원시 지표는 \tabref{tab:app-operationalization}에 정리하였다.}

\begin{table}[htbp]
\centering
\caption{\syncr{\chg{관측성과 불확실성, $k_O$}의 1차 조작화 \m{예시(참고용)}}}
\label{tab:app-operationalization}
\small
\syncr{\begin{tabularx}{\textwidth}{l l X l}
\toprule
\textbf{주 변수} & \textbf{1차 차원} & \textbf{원시 \m{지표(정지점)}} & \textbf{측정 출처} \\
\midrule
\multirow{5}{*}{$O_{\text{raw}}$} &
\multirow{3}{*}{$O_{\text{APM}}$} &
  $R_{\text{avail}}$: 가용시간/총시간 & APM 도구 로그 \\
& & $R_{\text{error}}$: 오류건수/총요청수 & APM 도구 로그 \\
& & $R_{\text{perf}}$: 목표응답/평균응답 & APM 도구 로그 \\
\cmidrule(l){2-4}
& \multirow{2}{*}{$O_{\text{SIEM}}$} &
  $E_{\text{high}}$: 미대응 고위험 이벤트 비율 & SIEM 도구 \\
& & $T_{\text{mttr}}$: 평균 \erf{탐지시간}{복구시간}(정규화) & SIEM 도구 \\
\midrule
\multirow{2}{*}{\jrev{$U_{\text{raw}}$}} &
  $U_{\text{hist}}$ & $\ln(P99/\jrev{\overline{T}})$ of 복구시간 & 장애 이력 DB \\
& $U_{\text{ext}}$ & \jrev{CVE CVSS 평균} & NVD/MITRE ATT\&CK \\
\midrule
$k_O$ & (단일) & 모니터링 커버리지율 & 자산 대장 \\
\bottomrule
\end{tabularx}}
\end{table}

\syncr{\chg{다만} 가중치 $w_i$, $v_i$, $\alpha$는 본 예시에서
등가중 또는 0.7로 임의 가정하였으며, \chg{산업별과 조직별} 최적값은 향후 연구에서
실증 검증되어야 한다.}

\syncr{\subsubsection{불확실성 \imod{$U$}{$U_{\text{raw}}$} 조작화 예시}}

\syncr{불확실성 \jrev{$U_{\text{raw}}$}는 과거 장애 이력 기반 $U_{\text{hist}}$와 외부 위협 지수 기반
$U_{\text{ext}}$의 두 차원으로 분해 가능하다\jrev{(두 하위 지표 모두 $[0, \infty)$
범위로 본 연구 모델의 $U_{\text{raw}}$ 정의와 일치한다)}.}

\syncr{\chg{과거 장애 기반 지표는 $U_{\text{hist}} = \ln\!\left(P99_{\text{recover}} / \overline{T}_{\text{recover}}\right) \jrev{\in [0, \infty)}$로, 복구시간 분포의 꼬리 두께를 로그 스케일로 정량화하며 장애 이력 DB에서 직접 산출한다.}}

\syncr{\chg{외부 위협 기반 지표는 $U_{\text{ext}} = \text{CVSS}_{\text{avg}} \jrev{\in [0, \infty)}$로, 최근 12개월 관련 CVE의 CVSS 점수 평균이며 NVD/MITRE ATT\&CK 공개 DB에서 산출한다.} \jrev{표준 CVSS v3.1은 $[0, 10]$이지만 조직 내부 위협 가중치 적용 시 10 초과 가능하다.}}

\syncr{\chg{두 차원의 합성은 \jrev{$U_{\text{raw}}$} $= \beta \cdot U_{\text{hist}} + (1-\beta) \cdot U_{\text{ext}}$이며, 본 예시에서 $\beta = 0.6$으로 가정한다.}}

\syncr{\chg{다만} 조직별 이력 DB 구축 수준이 상이하므로 $U_{\text{hist}}$
산출의 신뢰도는 조직마다 다를 수 있다. 신설 조직 또는 이력 부족 조직의
경우 $\beta$를 낮추거나 업종 평균을 사용하는 보정이 필요하다. \jrev{또한
$U_{\text{hist}}$와 $U_{\text{ext}}$는 각각 독립적으로 가우시안 또는 파레토
분포를 따를 수 있으며(조합 시 $2 \times 2 = 4$가지 분포 시나리오 가능),
본 연구 조건 C2/C3는 양 극단(전량 가우시안 / 전량 파레토)만 다룬다.
혼합 분포 시나리오의 실증은 향후 연구 과제이다.}}

\syncr{\subsubsection{관측성 가중치 $k_O$ 조작화 예시}}

\syncr{관측성 가중치 $k_O \in [0,1]$는 조직이 관측성에 얼마나 투자했는지를
나타내는 감도 계수로, 본 예시에서는 \textbf{모니터링 커버리지율}로 다음과 같이 단순화한다\chg{.}}

\syncr{$k_O = \frac{\text{감시 대상 자산 수}}{\text{전체 핵심 자산 수}}$}

\syncr{자산 대장에서 직접 산출되는 원시값이며, 감시 대상 자산은 APM/SIEM 도구에
등록된 엔드포인트 또는 프로세스를 기준으로 한다.}

\syncr{\chg{다만} 커버리지율은 ``감시의 폭''만을 측정하며, ``감시의
깊이''(로그 수집 상세도, 알람 정확도)는 별도 지표로 보완되어야 한다.
본 예시는 최소 구성만 제시한다.}

\syncr{\subsubsection{가상 기업 라이프사이클 적용 시나리오}}

\syncr{``(주)한빛정밀''(가상 기업명, 동명의 실제 기업과 무관)은
수도권 산업단지에 소재한 중견 제조업체(반도체 부품 정밀가공, 직원 420명,
연매출 1,200억 원 규모)로 가정하여 다음과 같이 시나리오를 작성한다. \nrev{ISO 22301:2019} 인증 준비 단계이며,
APM(Datadog)과 SIEM(Splunk)을 \chg{구축하여 운영한다}.}

\syncr{\chg{1단계인 평상시 진단(T-365일)에서는} 연간 BCMS 성숙도 평가 시점의 원시 지표를 다음과 같이
$R_{\text{avail}} = 0.994$, $R_{\text{error}} = 0.015$, $R_{\text{perf}} = 0.82$,
$E_{\text{high}} = 0.08$, $T_{\text{mttr}} = 0.65$, 커버리지율 $= 0.72$로 산출하였다.}

\syncr{합성 결과는 다음과 같다\chg{.} $O_{\text{APM}} = 0.933$, $O_{\text{SIEM}} = 0.785$,
$O_{\text{raw}} = 0.889$, $k_O = 0.72$, $U = 0.55$로 산출되어
$\eta \approx 1.05$ (Level 2 Repeatable) 판정된다. 모니터링 커버리지 확대 및
미대응 이벤트 축소가 개선 필요 영역으로 도출된다.}

\syncr{\chg{2단계인 평상시 개선 활동(T-365일부터 T-1일)에서는} 분기별 복구 훈련 도입, 모니터링 자산 확대(커버리지율 0.72 → 0.88),
SIEM 튜닝($E_{\text{high}}$ 0.08 → 0.03), 장애 이력 DB 정비를 수행하였다.
금년 재평가\chg{에서} $\eta \approx 0.82$\chg{를 달성하여} \textbf{\erf{Level 4 Managed}{Level 3 Defined}}로 \chg{향상되었다}.
월별 $\eta$ 추이를 경영진 대시보드에 정착한다.}

\syncr{\chg{3단계인 위기 발생 및 에스컬레이션(T+0h--T+6h)의} 가상 시나리오는 생산라인 제어 PLC 장애와 보조 전원계통 이상이 발생하였고
(인명 피해 없음, 설비 손상만 발생), 시간대별 $\eta$ 추이와 등급 변화는
\tabref{tab:app-lifecycle-eta}에 정리한다.}

\begin{table}[htbp]
\centering
\caption{\syncr{가상 기업 4단계 라이프사이클 $\eta$ \m{추이(참고 시나리오)}}}
\label{tab:app-lifecycle-eta}
\small
\syncr{\begin{tabularx}{\textwidth}{c l X c c}
\toprule
\textbf{시각} & \textbf{단계} & \textbf{상황} & \textbf{$\eta$} & \textbf{등급} \\
\midrule
$T{-}365$d & 1단계 진단 & 초기 \m{상태(Level 2)} & 1.05 & \erf{Yellow}{Orange} \\
$T{-}1$d & 2단계 개선 후 & 모니터링 확대 + SIEM \m{튜닝(\erf{Level 4}{Level 3})} & 0.82 & \erf{Green}{Yellow} \\
\midrule
$T{+}0$h & 3단계 정상 & 사건 직전 & 0.82 & \erf{Green}{Yellow} \\
$T{+}1$h & 3단계 탐지 & PLC 장애 탐지 & 1.15 & Orange \\
$T{+}2$h & 3단계 확산 & 전원계통 이상 확산 & 1.48 & Orange \\
$T{+}3$h & 3단계 정점 & 복구 불확실성 급등 & 1.62 & \textbf{Red} \\
$T{+}4$h & 3단계 전환 & 외부 지원 투입 & 1.45 & Orange \\
$T{+}6$h & 3단계 부분복구 & 주요 지표 안정화 & 1.08 & \erf{Yellow}{Orange} \\
\midrule
$T{+}72$h & 4단계 환류 & 사후 분석 후 재평가 & 0.84 & \erf{Green}{Yellow} \\
\bottomrule
\end{tabularx}}
\end{table}

\syncr{에스컬레이션\chg{은} $T{+}1$h(Orange)\chg{에} \chg{복구팀과 부서장} 자동 통지\chg{로 발동되어,}
$T{+}3$h(Red) 경영진 \chg{긴급 보고와 BCP 전면 가동을 거쳐,} $T{+}4$h 국가 위기경보 체계
연동 검토로 단계화되며 5.3.3.3(위기 시 정책 활용 방안)에서 상세 내용을 참조한다.}

\syncr{\chg{4단계인 복구와 환류(T+24h--T+72h)의} 사후 분석에서는 실제 복구 시간 $T_{\text{actual}} = 8.2$h vs $\DRTO$ 예측 7.6h
→ 상대 오차 $E_r = 7.3\%$ (수용 범위)이다. $O(t)$ 사각지대로 전원계통 센서
커버리지 부족을 식별하고, 성숙도 ``Managed → Optimized'' 이행 로드맵에
반영한다. 이번 사건을 $U_{\text{hist}}$ DB에 환류하여 내년도 진단의 입력으로 사용한다.}

\syncr{\chg{다만} 본 시나리오의 모든 수치는 중견 제조업 일반 특성에서
추정한 \textbf{조건적 예시}이며, 실제 기업의 실적 데이터가 아니다.
개별 조직의 적용 시 파라미터 재산출과 민감도 분석이 필수적이다.}

\syncr{\subsubsection{원시 지표별 측정 도구 매핑 예시}}

\syncr{\chg{앞의 관측성 $O_{\text{raw}}$, 불확실성 $U_{\text{raw}}$, 가중치 $k_O$ 조작화 예시에서 정의한} 원시 지표별 대표 측정 도구 매핑은
\tabref{tab:app-tool-mapping}에 정리하였다. 본 매핑은 \textbf{참고용}이며,
특정 도구의 추천이 아니다.}

\begin{table}[htbp]
\centering
\caption{\syncr{원시 지표별 측정 도구 매핑 \m{예시(참고용, 도구 권장 아님)}}}
\label{tab:app-tool-mapping}
\small
\syncr{\begin{tabularx}{\textwidth}{l l X}
\toprule
\textbf{원시 지표} & \textbf{측정 도구 유형} & \textbf{대표 \m{도구(참고)}} \\
\midrule
$R_{\text{avail}}$, $R_{\text{error}}$, $R_{\text{perf}}$ & APM & Datadog, New Relic, Prometheus+Grafana, Dynatrace \\
\addlinespace
$E_{\text{high}}$, $T_{\text{mttr}}$ & SIEM & Splunk, IBM QRadar, ArcSight, Elastic SIEM \\
\addlinespace
$U_{\text{hist}}$ & 장애 이력 DB & ServiceNow IRM, Jira Service Management, 자체 ITSM \\
\addlinespace
$U_{\text{ext}}$ & 위협 인텔리전스 & NVD, MITRE ATT\&CK, Recorded Future, CrowdStrike \\
\addlinespace
$k_O$ & 자산 대장/CMDB & ServiceNow CMDB, Device42, Lansweeper \\
\bottomrule
\end{tabularx}}
\end{table}

\nrev{ISO/IEC 25010:2011} \syncr{품질 모델의 Reliability(신뢰성),
Performance Efficiency(성능), Security(보안) 하위 속성이 각 원시 지표와
자연스럽게 대응되므로, 국제 표준 품질 지표 프레임워크와의 연계가 가능하다.}

\syncr{\chg{다만} 상기 도구는 예시이며, 조직의 기존 \chg{인프라와 예산, 규제}
환경에 따라 대체 선택이 필요하다.}

\syncr{\subsubsection{한계 종합 및 향후 실증 과제}}

\syncr{하위 변수 조작화와 시나리오는 참고용 예시이며, 다음 네 영역에서
후속 실증 연구가 필요하다.}

\syncr{\chg{첫째, 가중치 실증 튜닝이다.} $w_i$, $v_i$, $\alpha$, $\beta$를 \chg{산업별·조직 규모별}로 조정하는 다중 회귀 또는 최적화 연구가 필요하며,
본 예시는 등가중 또는 임의 고정값을 사용하였다.}

\syncr{\chg{둘째, 측정 도구 신뢰도이다.} 각 원시 지표의 측정 도구별 \chg{정확도와 편향}
검증이 필요하며, APM/SIEM 도구별 측정 알고리즘이 상이하므로
교차 타당도(cross-validity) 연구가 선결 과제이다.}

\syncr{\chg{셋째, 내용 타당도 재해석이다.} 본 연구의 \jrev{H9} 전문가 검증
(Section C 파라미터 적정성 6문항, S-CVI/Ave $= 0.985$)은 조작화의 내용
타당도의 간접 근거로 재해석 가능하나, 조작화 구성 자체에 대한 독립적
내용 타당도 검증은 별도 연구가 필요하다.}

\syncr{\chg{넷째, 수렴 및 예측 타당도이다.} 실제 조직 장애 이력과 $\eta$ 예측값의
상관 분석(종단 연구)을 통해 \chg{수렴, 판별, 예측} 타당도를 확보하는 실증 연구가
필요하다. 이는 본 연구 5.4.1 ``시뮬레이션 기반 검증의 한계''와 5.4.4
``시간 동역학의 부재''에서 제시한 향후 연구 방향과 연결된다.}

\syncr{종합적으로 하위 변수 조작화는 $\DRTO$ 프레임워크의 실무 적용 \textbf{방향성}을
제시하였으나, 본격 실무 배포 이전에는 상기 네 영역의 실증 연구가 선결되어야
한다.}



\FloatBarrier


% ===================================================================
% 5.3.4 산업별 적용 가이드라인 [구 5.3.2]
% ===================================================================
\subsection{\chg{산업별 적용 가이드라인}}
\label{subsec:practical-industry}

\begin{p5add}
\chg{본 항은 $\DRTO$ 모델의 산업별 적용 매트릭스를 제공한다. 먼저 P85.8 교차점에 기반한 2단계
프레임워크(Core C2 + Tail C3)로 일상 운영과 극단 시나리오의 분리 운용 원칙을 정립하고, 이어 6개 대표
산업(제조, 금융, 의료, IT 서비스, 물류, 소매)에 대한 $R_0$/$R_{\max}$/$\kappa$ 설정값과 적용 사례를
매트릭스 형태로 제시한다. 산업별 권장은 앞 항들의 회귀식 산출(\subsecref{subsec:practical-regression-tool})과
이중채널 자원 배분(\subsecref{subsec:practical-dualchannel-tool})을 각 산업의 분포 특성에 적용한 결과이다.}
\end{p5add}

\begingroup
  \renewcommand{\section}[1]{}
  \renewcommand{\subsection}[1]{\subsubsection{\syncr{#1}}}

  %% -------------------------------------------------------------------
\subsection{\vrev{Core(C2)와 Tail(C3)의 2단계 프레임워크}}
\label{subsec:two-stage-framework}
%% -------------------------------------------------------------------

앞서 논의한 교차점 분석과 이중채널 감쇠 결과에 기반하여,
일상 운영과 극단 시나리오를 분리하는 2단계(two-stage) $\DRTO$ 운영 프레임워크를
제안한다. \p{[}그림~\vref{fig:ch6-implementation-framework}\p{]}는 모니터링에서 에스컬레이션까지의
전체 운영 흐름을 도식화한 것이다.
Core(C2)와 Tail(C3)의 2단계 프레임워크에서는 2단계 분리의 원리와 Core/Tail 운영 전략을 중심으로 기술하며,
$\eta$ 기반 에스컬레이션 등급 체계와 구체적 대응 절차는
\subsecref{subsec:escalation-thresholds}에서 상세히 다룬다.

\begin{figure}[htbp]
\centering
\resizebox{\textwidth}{!}{%
\begin{tikzpicture}[
    stepbox/.style={draw, rounded corners=6pt, minimum width=36mm, minimum height=20mm,
                 align=center, font=\small, line width=0.7pt},
    detail/.style={draw, rounded corners=3pt, fill=white, minimum width=34mm,
                   align=left, font=\scriptsize, inner sep=4pt},
    conn/.style={-{Stealth[length=2.5mm]}, thick, blue!60},
    feedback/.style={-{Stealth[length=2mm]}, thick, black, dashed},
]

% 4단계 박스
\node[stepbox, fill=green!12]  (s1) at (0, 0)
  {\textbf{1. 모니터링}\\$O(t)$, $U(t)$ 수집\\APM/SIEM 연동};
\node[stepbox, fill=blue!12]   (s2) at (4.5, 0)
  {\textbf{2. 계산}\\$\DRTO$ 산출\\C2 또는 C3 적용};
\node[stepbox, fill=orange!12] (s3) at (9, 0)
  {\textbf{3. 의사결정}\\$\eta$ 평가\\Core vs Tail 판별};
\node[stepbox, fill=red!12]    (s4) at (13.5, 0)
  {\textbf{4. 에스컬레이션}\\$\eta$ 기반\\단계별 대응};

% 연결
\draw[conn] (s1) -- (s2);
\draw[conn] (s2) -- (s3);
\draw[conn] (s3) -- (s4);

% 세부 내용 (피드백 루프보다 먼저 배치)
\node[detail] (d1) at (0, -2.5) {
  $\bullet$ 관측성: APM 지표 $\to$ $O \in [0,1]$\\
  $\bullet$ 불확실성: 장애 이력 $\to$ $U$\\
  $\bullet$ 30분 주기 갱신\\
  $\bullet$ 이상 탐지 시 즉시 갱신
};
\node[detail] (d2) at (4.5, -2.5) {
  $\bullet$ Core ($\leq$ P85.8): C2 적용\\
  $\bullet$ Tail ($>$ P85.8): C3 적용\\
  $\bullet$ $\DRTO = R_0 + E_{O,bnd} + E_{U,bnd}$\\
  $\bullet$ $\eta = \DRTO / R_0$ 산출
};
\node[detail] (d3) at (9, -2.5) {
  $\bullet$ $\eta < 0.8$: 안정 구간\\
  $\bullet$ $0.8 \leq \eta < 1.0$: 주의\\
  $\bullet$ $1.0 \leq \eta < 1.5$: 경고\\
  $\bullet$ $\eta \geq 1.5$: 위기
};
\node[detail] (d4) at (13.5, -2.5) {
  $\bullet$ 단계별 보고 대상 결정\\
  $\bullet$ 추가 자원 배정\\
  $\bullet$ 비상 계획 가동 판단\\
  $\bullet$ 사후 분석 및 교정
};

% 세부 연결
\draw[dashed, black, line width=0.8pt] (s1.south) -- (d1.north);
\draw[dashed, black, line width=0.8pt] (s2.south) -- (d2.north);
\draw[dashed, black, line width=0.8pt] (s3.south) -- (d3.north);
\draw[dashed, black, line width=0.8pt] (s4.south) -- (d4.north);

% 피드백 루프 (상단 우회: s4 → 위 → s1)
\draw[feedback] (s4.north) -- ++(0, 1.0) -| (s1.north)
    node[pos=0.5, above, font=\scriptsize, text=black] {파라미터 교정 피드백};

% 단계 레이블 (피드백 루프 위에 배치)
\node[font=\footnotesize\bfseries, text=black] at (0, 2.8)  {실시간 입력};
\node[font=\footnotesize\bfseries, text=black]  at (4.5, 2.8) {모델 적용};
\node[font=\footnotesize\bfseries, text=black] at (9, 2.8) {상태 판별};
\node[font=\footnotesize\bfseries, text=black]   at (13.5, 2.8) {대응 실행};

\end{tikzpicture}
}% end resizebox
\caption{모니터링에서 계산, 의사결정, 에스컬레이션으로 이어지는 $\DRTO$ 구현 프레임워크}
\label{fig:ch6-implementation-framework}
\end{figure}

2단계 프레임워크의 핵심 원리는 \p{[}표~\vref{tab:two-stage-summary}\p{]}와 같다.

\begin{table}[htbp]
\centering
\caption{Core(C2)와 Tail(C3)의 2단계 프레임워크 요약}
\label{tab:two-stage-summary}
\small
\begin{tabularx}{\textwidth}{l X X}
\toprule
\textbf{구분} & \textbf{1단계 --- Core \m{영역(C2 기반)}} & \textbf{2단계 --- Tail \m{영역(C3 기반)}} \\
\midrule
\jrev{적용 \m{영역(P85.8 분할)}}     & P85.8 \m{이하(일상적 운영)}             & P85.8 \m{초과(극단 시나리오)} \\
\addlinespace
\jrev{환경 분포}  & 가우시안 ($U \sim N(3,1)$)           & 파레토 ($U \sim 6 \cdot \mathrm{Pa}(3)$) \\
\addlinespace
\jrev{환경 전체 평균 $\bar{\eta}$}  & $\bar{\eta}_{\text{C2}} \approx 1.68$ & $\bar{\eta}_{\text{C3}} \approx 1.51$ \\
\addlinespace
\jrev{C3 내 분할 영역 $k_O$ 계수}    & \jrev{Core: $-0.798$}                             & \jrev{Tail: $-0.415$ ($0.52$배 감쇠)} \\
\addlinespace
CVaR$_{99}$ \jrev{(환경 전체)}   & $3.180$                              & $3.944$ ($+24.0\%$) \\
\addlinespace
\jrev{환경 전체 $R^2$ (9변수)} & $0.998$                              & $0.858$ \\
\addlinespace
운영 목표      & 효율적 \m{복구(관측성 효과 극대화)}       & 보수적 \m{대응(꼬리 위험 대비)} \\
\addlinespace
자원 배분      & 정상 복구 절차                        & 추가 자원 투입, 비상 계획 연계 \\
\bottomrule
\end{tabularx}
\end{table}

\jrev{[표 5-14]는 Core/Tail의 두 가지 의미를 함께 다룬다. \chg{첫째는 P85.8 기준 \textit{분할 영역}(적용 영역 행)이고, 둘째는 분할 영역에 적용되는 \textit{환경 모델 C2/C3}(환경 모델, 환경 평균, 환경 $R^2$ 행)이며, 셋째는 C3 조건 \textit{내부}의 분할 회귀 결과($k_O$ 계수 행, 4.10절 dual\_channel 결과)이다.} 각 지표의 산출 분석 대상이 다르므로 행 단위로 의미를 구분하여 해석해야 한다.}

\textbf{1단계(Core, C2 기반)}는 P85.8 이하의 일상적 운영 상황에서 적용되며,
가우시안 불확실성 분포(C2 조건)에 기반한 효율적 복구를 추구한다.
\jrev{C2 조건 전체} 평균 $\bar{\eta} \approx 1.68$이며,
\jrev{C3 조건의 분할 회귀 결과 Core 영역의 $k_O$ 계수는 $-0.798$로} 관측성의 효과가 충분히 발현된다.
\textbf{2단계(Tail, C3 기반)}는 P85.8 초과의 극단 시나리오에서 활성화되며,
파레토 불확실성 분포(C3 조건)의 특성을 반영한다.
\jrev{C3 분할 회귀 결과 Tail 영역의 $k_O$ 계수가 $-0.415$로 $0.52$배 감쇠되고},
CVaR$_{99}$가 C2 대비 $+24.0\%$ 높아져, 불확실성 관리 중심의 추가 자원
투입이 정당화된다.

이러한 2단계 분리는 앞서 규명된 P85.8 분할 기준점의 이론적 근거에
기반하며, 조직이 일상 운영에서는 효율성을, 극단 상황에서는
안정성을 각각 최적화할 수 있는 구조적 기반을 제공한다.


  
%% -------------------------------------------------------------------
\subsection{산업별 권장 파라미터}
\label{subsec:industry-guidelines}
%% -------------------------------------------------------------------

$\DRTO$ 프레임워크의 파라미터는 산업별 특성에 따라 교정(calibration)이
필요하다. \chref{ch:results}의 시뮬레이션 결과, 전문가 의견, 그리고 연구모형의 산업
시나리오 설계를 종합하여, 6개 주요 산업에 대한 권장 파라미터 설정을
\p{[}표~\vref{tab:ch6-industry-params}\p{]}에 제시한다.

\begin{table}[htbp]
\centering
\caption{6개 산업별 $\DRTO$ 권장 파라미터}
\label{tab:ch6-industry-params}
\small
\begin{tabularx}{\textwidth}{l c c c c X}
\toprule
\textbf{산업} & \textbf{$R_0$ (h)} & \textbf{$R_{\max}$ (h)} & \textbf{$O_{\text{raw}}$ 범위} & \textbf{권장 조건} & \textbf{적용 근거} \\
\midrule
제조업       & 4.0 & 16.0 & 0.50--0.70 & C2
  & 생산 설비 고장은 가우시안 분포 적합, 생산 계획 연동으로 $R_0$ 엄격 관리 필요 \\
\addlinespace
금융         & 1.0 & 4.0  & 0.70--0.85 & C3
  & 규제 준수 요건(RTO $\leq$ 2h) 엄격, 시장 변동$\cdot$사이버 공격의 파레토 특성 반영 필수 \\
\addlinespace
의료         & 2.0 & 8.0  & 0.60--0.75 & C3
  & 환자 안전이 최우선, 의료기기 장애, 감염 확산 시나리오의 파레토 위험 대비 \\
\addlinespace
IT 서비스    & 0.5 & 2.0  & 0.75--0.90 & C2
  & SLA 기반 복구, 인프라 장애의 대다수가 가우시안 근사 가능, 높은 관측성 인프라 \\
\addlinespace
물류         & 6.0 & 24.0 & 0.45--0.60 & C2
  & 공급망 연속성, 계절 변동$\cdot$물류 지연은 정규분포 근사, 가장 긴 허용 RTO \\
\addlinespace
소매         & 3.0 & 12.0 & 0.50--0.65 & C2
  & 매출 손실 최소화, 일상적 IT 장애 중심, POS$\cdot$결제 시스템 관측성 활용 \\
\bottomrule
\end{tabularx}
\end{table}

\chg{[표 5-15]의 공통 파라미터는 다음과 같다. 곡률 매개변수는 $\kappa = 1.2$로 \jrev{\citet{lee2026drto}}에서 확정된
매개변수이며, 본 연구는 \frev{분포 환경 적합성} 종합 명제로 \frev{운영 적합성을 시사한다}(\secref{sec:ch4_hypothesis_summary}
참조). 관측성 가중치는 권장 범위 $k_O \in [0.4, 0.6]$이고, C2는 가우시안 불확실성 분포, C3는 파레토
불확실성 분포이며, $R_0 / R_{\max}$은 연구모형 기본 설정으로 1:4 비율을 유지한다. $O_{\text{raw}}$ 범위는 해당 산업의
전형적 모니터링 성숙도를 반영한다.}

\p{[}그림~\vref{fig:ch6-industry-summary}\p{]}는 6개 산업의 핵심 파라미터와
권장 조건을 종합적으로 시각화한 것이다.

\begin{figure}[htbp]
\centering
\resizebox{\textwidth}{!}{%
\begin{tikzpicture}[
    indbox/.style={draw, rounded corners=5pt, minimum width=30mm, minimum height=22mm,
                   align=center, font=\small, line width=0.6pt},
    modelC2/.style={draw, rounded corners=3pt, fill=blue!10, minimum width=12mm,
                    minimum height=6mm, font=\scriptsize\bfseries},
    modelC3/.style={draw, rounded corners=3pt, fill=red!10, minimum width=12mm,
                    minimum height=6mm, font=\scriptsize\bfseries},
    conn/.style={-{Stealth[length=2mm]}, thick, black},
]

% 산업 박스
\node[indbox, fill=green!8]  (mfg) at (0, 0)
  {\textbf{제조업}\\$R_0 = 4.0$h\\$O = 0.50$--$0.70$};
\node[indbox, fill=blue!8]   (fin) at (3.2, 0)
  {\textbf{금융}\\$R_0 = 1.0$h\\$O = 0.70$--$0.85$};
\node[indbox, fill=red!8]    (hc)  at (6.4, 0)
  {\textbf{의료}\\$R_0 = 2.0$h\\$O = 0.60$--$0.75$};
\node[indbox, fill=purple!8] (it)  at (9.6, 0)
  {\textbf{IT 서비스}\\$R_0 = 0.5$h\\$O = 0.75$--$0.90$};
\node[indbox, fill=orange!8] (log) at (12.8, 0)
  {\textbf{물류}\\$R_0 = 6.0$h\\$O = 0.45$--$0.60$};
\node[indbox, fill=yellow!10](ret) at (16.0, 0)
  {\textbf{소매}\\$R_0 = 3.0$h\\$O = 0.50$--$0.65$};

% 모델 레이블
\node[modelC2] at (0, -1.7) {C2};
\node[modelC3] at (3.2, -1.7) {C3};
\node[modelC3] at (6.4, -1.7) {C3};
\node[modelC2] at (9.6, -1.7) {C2};
\node[modelC2] at (12.8, -1.7) {C2};
\node[modelC2] at (16.0, -1.7) {C2};

% 연결
\draw[dashed, black] (mfg.south) -- (0, -1.4);
\draw[dashed, black] (fin.south) -- (3.2, -1.4);
\draw[dashed, black] (hc.south)  -- (6.4, -1.4);
\draw[dashed, black] (it.south)  -- (9.6, -1.4);
\draw[dashed, black] (log.south) -- (12.8, -1.4);
\draw[dashed, black] (ret.south) -- (16.0, -1.4);

% R₀ 크기 축 (상단)
\draw[thick, -{Stealth[length=2.5mm]}, black] (-1.2, 2.0) -- (17.2, 2.0)
    node[right, font=\scriptsize, black] {$R_0$ 증가};
\node[font=\scriptsize, black, above] at (3.2, 2.0) {짧은 RTO};
\node[font=\scriptsize, black, above] at (12.8, 2.0) {긴 RTO};

% 범례
\node[modelC2] at (5.5, -2.8) {C2};
\node[font=\scriptsize, anchor=west] at (6.2, -2.8) {가우시안 (일상적 장애)};
\node[modelC3] at (10.5, -2.8) {C3};
\node[font=\scriptsize, anchor=west] at (11.2, -2.8) {파레토 (꼬리 위험 필수)};

\end{tikzpicture}
}% end resizebox
\caption{6개 산업 $\DRTO$ 적용 요약 다이어그램}
\label{fig:ch6-industry-summary}
\end{figure}

산업별 조건 선택의 핵심 기준은 \textbf{파레토 위험의 존재 여부}이다.
금융과 의료는 사이버 공격, 시장 급변, 감염 확산 등 극단적 사건이
복구 시간을 급격히 연장할 수 있으므로 C3(파레토)를 권장한다.
C3의 CVaR$_{99}$가 C2 대비 $+24.0\%$ 높은 점은 이러한 산업에서
꼬리 위험을 사전에 반영해야 하는 정량적 근거가 된다.
반면, 제조업, IT 서비스, 물류, 소매는 장애의 대다수가 가우시안
분포로 근사 가능하므로 C2가 효율적이며, 필요 시 극단 시나리오에 한하여
C3를 보완적으로 적용할 수 있다.

분포 유형의 판단 오류는 복구 자원의 배정에서 두 가지 상반된 위험으로 이어진다.
가우시안 환경(C2)에 파레토 수준(C3)의 자원을 배정하는 과다예비(over-provisioning)는
불확실성을 과대 추정한 결과로, 유휴 대기 인력이나 미사용 예비 시스템, 과도한 예비
부품과 같은 자원이 상시 유지되어 비용 효율이 저하되나 서비스 수준 협약(SLA) 위반
위험은 낮다. 예를 들어 정상적 IT 장애 복구에 대해 랜섬웨어 대응 수준의 인력과 예비
시스템을 상시 유지하면 유휴 자원 비용이 불필요하게 증가한다. 반대로 파레토
환경(C3)에 가우시안 수준(C2)의 자원만 배정하는 과소예비(under-provisioning)는
불확실성을 과소 추정한 결과로, 비용은 절감되나 극단적 사건 발생 시 SLA 위반과
업무 중단 장기화 위험이 급증한다. 공급망 복합 장애 가능성이 있는 물류 기업이 단순
장비 교체 수준의 예비만 확보하면, P99 수준의 극단적 지연(약 $13.2$시간 이상)이
발생할 때 복구 역량이 부족하여 업무 연속성이 심각하게 위협받는다. 따라서 산업별
조건 선택은 단순한 분포 가정의 문제가 아니라 자원 배정의 비용과 위험을 균형 짓는
의사결정이며, 파레토 위험이 존재하는 산업에서 C3를 권장하는 근거가 여기에 있다.

각 산업에 대한 세부 적용 지침은 다음과 같다.

\chg{제조업의 경우,}
생산 설비 고장에 의한 업무 중단이 주요 위험이며,
관측성은 PLC(Programmable Logic Controller), SCADA, MES(Manufacturing
Execution System) 등의 기존 인프라를 활용한다.
$R_0 = 4.0$h는 일반적인 교대제 생산 라인의 복구 목표를 반영하며,
$O = 0.50$--$0.70$은 중간 수준의 모니터링 성숙도에 해당한다.

\chg{금융의 경우,}
금융감독 당국의 IT 복구 시간 규제(예: RTO $\leq$ 2시간)가 엄격하며,
\chg{사이버 공격, 시장 급변, 결제 시스템 장애}의 파레토 특성을 고려하여
C3를 권장한다. $O = 0.70$--$0.85$의 높은 관측성은 실시간 거래
모니터링 시스템(FDS 등)을 반영한다.

\chg{의료의 경우,}
환자 안전이 최우선이며, 의료기기 장애, 전자 의무기록(EMR) 시스템 중단,
감염병 확산 시나리오에서 복구 지연이 생명에 직결된다.
$R_0 = 2.0$h는 응급 의료 시스템의 전형적 복구 목표이며,
C3 적용으로 극단 시나리오에 대한 사전 대비를 강화한다.

\chg{IT 서비스의 경우,}
SLA(Service Level Agreement) 기반의 엄격한 복구 목표를 가지며,
$R_0 = 0.5$h는 클라우드/SaaS 서비스의 전형적 RTO이다.
$O = 0.75$--$0.90$의 높은 관측성은 APM, 로그 분석, 분산 추적(distributed
tracing) 등 성숙한 모니터링 인프라를 반영한다.

\chg{물류의 경우,}
공급망의 연속성이 핵심이며, 계절 변동, 운송 지연, 창고 시스템 장애가
주요 위험이다. $R_0 = 6.0$h는 본 연구의 기본 설정과 동일하며,
$O = 0.45$--$0.60$은 물류 산업의 상대적으로 낮은 IT 관측성 성숙도를 반영한다.

\chg{소매의 경우,}
POS(Point of Sale) 시스템, 재고 관리, 전자상거래 플랫폼의 가용성이
매출에 직접 영향을 미친다. $R_0 = 3.0$h, $O = 0.50$--$0.65$는
중간 규모 소매 환경의 전형적 특성을 반영하며, C2 기반의 효율적 복구를 권장한다.


\chg{이상의 산업별 권장 파라미터를 회귀식에 대입한 구체적 $\DRTO$ 산출 예시는 회귀식 기반 산출 도구를
다루는 \subsecref{subsec:practical-regression-tool}에 제시한다.}



\endgroup
\FloatBarrier
      % 5.3 실무적 기여
% =====================================================================
% 5.4 한계점 — 5단계 5-2 wrapper (sec5.4_limitations.tex + 5.1.3 흡수)
% 본문 텍스트 무수정. 원본 sec5.4_limitations.tex 그대로 + 5.1.3 본문 append.
% =====================================================================



%% ===================================================================
\section{\vrev{한계점}}
\label{sec:limitations}
%% ===================================================================

본 연구는 이론적 모델 수립과 시뮬레이션 기반 검증에 초점을 맞추고 있으며,
다음의 한계를 지닌다. \p{[}표~\vref{tab:ch6-limitations}\p{]}은 각 한계와
대응하는 향후 연구 방향을 종합한다.

\begin{table}[htbp]
\centering
\caption{\revdel{한계점 및 향후 연구 방향}\,\revadd{한계점 종합}}
\label{tab:ch6-limitations}
\small
\begin{tabularx}{\textwidth}{c l X X}
\toprule
\textbf{\#} & \textbf{한계 영역} & \textbf{한계 내용} & \textbf{향후 연구 방향} \\
\midrule
1 & 시뮬레이션 기반 검증 &
  $n = 10{,}000$ MC 시뮬레이션과 $10^8$회 다중 시드 검증은 통계적으로 강건하나,
  실제 조직 환경에서의 현장 실험(field test)을 수반하지 않음 &
  제조$\cdot$금융$\cdot$IT 서비스 등 산업별 파일럿 현장 검증, 실제 장애 이벤트에서의
  $\eta$ 예측 정확도 평가 \\
\addlinespace
2 & 고정 $R_{\max}/R_0$ 비율 &
  $R_0 = 6.0$h, $R_{\max} = 24.0$h ($1:4$ 비율) 고정,
  산업$\cdot$조직별 비율 변동 시 모델 행태 미탐색 &
  $R_{\max}/R_0$ 비율을 $[2, 8]$ 범위에서 변동시키는 민감도 분석,
  동적 $R_0(t)$ 모형으로의 확장 \\
\addlinespace
3 & 분포 가정 &
  C2: $U \sim N(3,1)$, C3: $6 \cdot \mathrm{Pareto}(\alpha=3)$ 고정,
  실제 불확실성은 다양한 꼬리 두께를 가질 수 있음 &
  \jrev{일반화 파레토 분포(GPD)와 GEV 분포 도입, 가우시안-파레토 혼합 분포 모형 개발, $\alpha$ 변동에 따른 민감도 분석(P85.8 교차점 이동 및 감쇠 비율 변화)} \\
\addlinespace
4 & 시간 동역학 부재 &
  현행 모델은 시점별 단면(cross-sectional) 분석,
  시간에 따른 $O(t)$, $U(t)$의 자기상관, 추세, 계절성 미반영 &
  \jrev{시계열 모형(ARIMA, GARCH) 통합, 베이지안 온라인 학습 기반 자기 교정 메커니즘 구축, 복구 과정 중 실시간 $\DRTO$ 갱신 프로토콜 수립} \\
\addlinespace
\tmp{5} & 심층 불확실성 미반영 &
  현행 모델은 분포 유형(가우시안/파레토)이 사전에 특정 가능하다는 전제,
  모델$\cdot$시나리오$\cdot$확률 모두 불확실한 심층 불확실성(deep uncertainty) 상황은
  적용 범위를 벗어남 &
  로버스트 최적화, 시나리오 발견,
  info-gap 의사결정 이론 등 분포 비의존적 접근법과의 통합 탐색 \\
\addlinespace
\tmp{6} & 국가 위기경보 체계와의 연계 부재 &
  $\eta$ 기반 에스컬레이션은 정량적 지표에 의해 등급이 자동 결정되나,
  국가 위기경보(관심--주의--경계--심각)의 질적 판단 체계와의
  공식적 연계 메커니즘이 마련되지 않음 &
  국가 위기경보 4단계와 $\eta$ 임계값의 매핑 프레임워크 개발,
  질적 판단과 정량적 지표의 상호 보완적 의사결정 모형 구축 \\
\addlinespace
\jrev{7} & \jrev{\frev{분포 환경 적합성의 간접 시사}} &
  \jrev{본 연구는 \chg{\citet{lee2026drto}}의 $\kappa^{*}=1.2$를 C2/C3 조건에 적용한 후 9개 가설 PASS의 통합 해석으로 \frev{운영 적합성을 시사하였으나}, 직접적 격자 탐색은 4기준 종합점수 $F(\kappa)$가 C1 특화 지표라 C2/C3에 부적합} &
  \jrev{환경 무관 다기준 종합점수 재정의, C2/C3 조건에 적합한 새로운 4기준 도출 후 직접 격자 탐색 수행} \\
\bottomrule
\end{tabularx}
\end{table}


%% -------------------------------------------------------------------
\subsection{\vrev{시뮬레이션 기반 검증의 한계}}
\label{subsec:lim-simulation}
%% -------------------------------------------------------------------

$n = 10{,}000$ 몬테카를로 시뮬레이션과 $10^8$회 다중 시드 검증은
통계적으로 강건한 결과를 제공하나, 실제 조직 환경에서의 현장 실험을
수반하지 않았다. 시뮬레이션에서 확인된 효과 크기
(Cohen's $d = -1.180$, $+1.871$)와 이중채널 감쇠($0.52$배)가
실제 복구 과정에서도 동일한 크기로 발현되는지는
추가적인 실증 연구를 통해 확인되어야 한다.
특히 관측성 지표 $O(t)$의 실시간 측정 정확도, 불확실성 분포의
실제 형태, 그리고 조직의 복구 역량 변동이 시뮬레이션 가정과
어떻게 괴리되는지를 현장 데이터로 검증하는 것이 필요하다.
$O_{\text{raw}}$의 조작적 정의로서, 장애 탐지율, 로그 가시성, 복구 진척 추적 수준 등의
하위 지표를 리커트 척도로 평가하여 $[0, 1]$로 정규화하는 방법을 예시할 수 있으나,
측정 지표의 가중치 산정과 \chg{산업별 및 조직별} 특성을 반영한 측정 도구의 정밀화는
후속 연구에서 수행되어야 한다.


%% -------------------------------------------------------------------
\subsection{\vrev{고정 $R_{\max}/R_0$ 비율의 한계}}
\label{subsec:lim-ratio}
%% -------------------------------------------------------------------

$R_{\max}/R_0$ 비율을 모든 산업에서 4로 고정하여 비율 변동의 영향을 탐색하지
않은 점은 본 연구의 한계이다. 본 연구에서는 $R_0 = 6.0$h, $R_{\max} = 24.0$h로 $R_{\max}/R_0 = 4$의
고정 비율을 사용하였다. 이 비율은 물류 산업의 전형적 설정(\p{[}표~\vref{tab:ch6-industry-params}\p{]})과
일치하나, 금융($R_{\max}/R_0 = 4$), IT 서비스($R_{\max}/R_0 = 4$) 등
모든 산업에서 동일 비율이 적용되었다. 실무에서는 \chg{산업과 조직, 규제} 환경에 따라
이 비율이 $2$에서 $8$ 이상까지 변동할 수 있으며, 비율 변동이 시그모이드의
작동 영역($z$ 범위)과 이중채널 감쇠 비율에 미치는 영향은
체계적으로 탐색되지 않았다.

향후 연구에서는 $R_{\max}/R_0$ 비율을 독립 변수로 설정한 민감도 분석을 통해,
비율 변동에 대한 모델의 강건성을 검증하고, 산업별 최적 비율 범위를
도출하는 것이 필요하다. 나아가 $R_0(t)$ 자체를 시간 변동 함수로 확장하여,
BIA 갱신 시마다 기준선이 적응적으로 조정되는 동적 기준선 모형을 개발할 수 있다.


%% -------------------------------------------------------------------
\subsection{\vrev{분포 가정의 한계}}
\label{subsec:lim-distribution}
%% -------------------------------------------------------------------

불확실성 분포를 가우시안과 파레토($\alpha = 3$)의 두 형태로 고정한 점은 본 연구의
또 다른 한계이다. C2 조건에서의 불확실성 분포 $U \sim N(3,1)$과 C3 조건에서의
$6 \cdot \mathrm{Pareto}(\alpha = 3)$은 각각 경꼬리와 중꼬리의
대표적 형태를 포착하기 위한 선택이다. 그러나 실제 재난 복구 과정에서의
불확실성은 $\alpha < 3$(더 무거운 중꼬리)이나 $\alpha > 3$(더 가벼운 경꼬리)의
다양한 특성을 보일 수 있다.

$\alpha$ 값의 변동에 따른 핵심 결과의 민감도\chg{(}P85.8 교차점의 이동,
CVaR 격차의 변화, 이중채널 감쇠 비율\chg{)}는 향후 연구에서 체계적으로 탐색되어야 한다.
일반화 파레토 분포(GPD)나 극단값 이론(EVT)의 GEV 분포를
불확실성 모형에 도입하여, 다양한 꼬리 두께에 대한 모델의 적응성을
확보하는 것이 바람직하다.

\jrev{또한 본 연구의 \textbf{\frev{분포 환경 적합성} 종합 명제}(4.13절)는 \chg{\citet{lee2026drto}}의
$\kappa^{*}=1.2$의 C2/C3 조건 적용성을 9개 가설(\jrev{H2--H10}) PASS의 통합 해석으로
\frev{시사한다}. 이는 직접적 격자 탐색의 대안으로서 학술적 정직을 우선한 접근이지만,
\chg{\citet{lee2026drto}}의 4기준 종합점수 $F(\kappa) = f_1 \cdot f_2 \cdot f_3 \cdot f_4$
(\chg{실무 유의성, 비포화, 효과 크기, 설명력})가 C1 조건 기반으로 정의되어
C2/C3 환경에서 직접 격자 탐색에 부적합하다는 \textbf{지표 일반화의 한계}를
내포한다. 환경 무관 다기준 종합점수의 재정의 및 C2/C3 환경에서의
직접 격자 탐색은 본 연구 범위를 벗어나는 후속 연구 과제이다.}


%% -------------------------------------------------------------------
\subsection{\vrev{시간 동역학의 부재}}
\label{subsec:lim-temporal}
%% -------------------------------------------------------------------

현행 $\DRTO$ 모델은 특정 시점의 $O(t)$와 $U(t)$를 입력으로 받아
해당 시점의 $\DRTO$를 산출하는 단면(cross-sectional) 모형이다.
시간에 따른 입력 변수의 자기상관(autocorrelation), 추세(trend),
계절성(seasonality)은 모형에 내생적으로 반영되지 않는다.

실제 복구 과정에서는 $O(t)$가 복구 진행에 따라 점진적으로 증가하고,
$U(t)$는 초기에 높았다가 정보 획득에 따라 감소하는 시간적 패턴을 보인다.
이러한 동역학을 포착하기 위해, 베이지안 온라인 학습(Bayesian online learning)을
적용하여 장애 이벤트 관측 시마다 매개변수를 적응적으로 갱신하는
자기 교정(self-calibrating) 메커니즘의 구축이 향후 연구 방향이다.
ARIMA, GARCH 등 시계열 모형~\citep{durbin2012timeseries,hyndman2018forecasting}과의 통합을 통해
$\DRTO$ 모형에 시간 차원을 추가하는 확장도 고려할 수 있다.


%% -------------------------------------------------------------------
\subsection{\vrev{심층 불확실성의 미반영}}
\label{subsec:lim-deep-uncertainty}
%% -------------------------------------------------------------------

심층 불확실성을 반영하지 못하는 점은 현행 모델의 구조적 한계이다.
현행 $\DRTO$ 모델은 불확실성의 확률분포 유형이 사전에 특정 가능하다는 전제 하에
설계되었다. C2 조건에서는 가우시안, C3 조건에서는 파레토 분포를 채택하여
각각 경꼬리와 중꼬리 특성을 포착한다. 그러나 팬데믹, 기후변화에 의한 복합재난,
전례 없는 대규모 사이버 공격 등 \textbf{심층 불확실성}(deep uncertainty) 상황에서는
분포의 형태는 물론 모델 구조와 시나리오 자체를 사전에 특정할 수 없다.

이는 \subsecref{subsec:uncertainty-risk}에서 정리한 불확실성 분류 체계의
네 번째 유형\chg{(}\chg{모델, 시나리오, 확률} 모두 불확실한 심층 불확실성\chg{)}에 해당하며,
현 프레임워크의 구조적 적용 한계를 구성한다.
나이트적(Knightian) 불확실성은 파레토 분포로 보수적 근사가 가능하나,
심층 불확실성은 분포 선택 자체가 불가능하므로 확률분포 기반 모형의 범위를 벗어난다.

향후 연구에서는 로버스트 최적화(robust optimization),
시나리오 발견(scenario discovery),
info-gap 의사결정 이론 등 분포 비의존적(distribution-free) 접근법을
$\DRTO$ 프레임워크와 통합하여,
심층 불확실성 상황에서도 복구시간의 동적 조정이 가능한 확장 모형을 탐색할 필요가 있다.


%% -------------------------------------------------------------------
\subsection{\vrev{국가 위기경보 체계와의 연계}}
\label{subsec:lim-national-alert}
%% -------------------------------------------------------------------

본 연구의 $\eta$ 기반 에스컬레이션은 국가 위기경보 체계와의 공식적 연계
메커니즘을 아직 갖추지 못하였다. \subsecref{subsec:escalation-thresholds}에서 논의한 바와 같이,
현행 국가 위기경보 체계(관심--주의--경계--심각)는
재난 및 안전관리 기본법에 근거한 \textbf{질적 판단}에 기반하여
전문가 합의와 상황 평가를 통해 등급이 결정된다.
반면 본 연구의 $\eta$ 기반 에스컬레이션은
$\DRTO / R_0$라는 \textbf{정량적 지표}에 의해 등급이 자동으로 산출되므로,
두 체계는 상호 보완적 관계를 형성할 잠재력을 지닌다.

구체적으로, 국가 위기경보가 발령된 상황에서 $\eta$ 값을 병행 산출하면
질적 판단의 객관적 근거를 제공할 수 있으며,
역으로 $\eta$의 급격한 상승이 감지될 때 위기경보 상향 검토의
조기 경고(early warning) 신호로 활용할 수 있다.
향후 연구에서는 국가 위기경보 4단계와 $\eta$ 임계값 간의
체계적 매핑 프레임워크를 개발하고,
질적 판단과 정량적 지표가 결합된 하이브리드 의사결정 모형을 구축하여
국가 수준의 재난 대응 체계에 $\DRTO$의 정량적 기여 가능성을
실증적으로 검증하는 것이 필요하다.


%% -------------------------------------------------------------------
\syncr{\subsection{하위 변수 조작화의 단순화}}
\syncr{\label{subsec:lim-operationalization}}
%% -------------------------------------------------------------------

\syncr{본 연구는 관측성($O_{\text{raw}}$, $k_O$)과 불확실성($U$)의
제1계층 하위 지표까지의 조작화 \textbf{예시}를 \subsecref{app:operationalization}에
참고용으로 제시하였으나, 가중치($w_i$, $\alpha$, $\beta$)의 산업별 튜닝,
원시 지표의 측정 도구별 \chg{정확도와 편향} 검증, 현장 데이터 기반 \chg{수렴 및 예측}
타당도 확보는 본 연구의 범위를 벗어나는 후속 연구 과제이다. 제시된
조작화 예시와 가상 기업 라이프사이클은 실무 적용의 방향성을 보여주기
위한 참고용일 뿐 검증된 방법론이 아니며, 심리측정학적 엄밀성을 갖춘
측정 도구 개발 및 종단 연구를 통한 예측 타당도 확보가 선결되어야 한다.}

\vspace{1em}

\fpara{본 연구의
가장 중요한 학술 한계는 \emph{$\kappa = 1.2$의 환경 한정성}이다. \chg{\citet{lee2026drto}}의 도출값
$\kappa = 1.2$는 C1 조건(관측성 단독, 가우시안 베이스라인) 격자 탐색에서 도출되었으며,
C2/C3 환경에서의 적합성은 4.13절의 H2--H10 PASS 통합 해석으로 \emph{운영 적합성을
시사}하나 \emph{환경별 최적성을 입증}하지는 못한다. 따라서 본 연구는 $\kappa = 1.2$를
BCM 일반 운영의 \emph{\chg{잠정적이고 보수적인} 실무 기본값}(provisional and conservative operational
default)으로 채택하며, \emph{환경 무관 보편 상수성을 주장하지 않는다}. 환경 무관 $\kappa$의
\emph{존재 여부}에 대한 다방법 탐색은 후속 연구 논문의 핵심 의제로 분리된다
(5.5.3절 향후 연구 참조).}

본 장에서 제시한 \chg{이론적, 방법론적, 실무적} 시사점과 한계는
$\DRTO$ 프레임워크의 현재 위상을 명확히 규정하는 동시에,
후속 연구의 구체적 방향을 제시한다. 특히 2단계 프레임워크(Core C2 + Tail C3),
$\eta$ 기반 에스컬레이션 체계, 6개 산업 적용 가이드라인은
$\DRTO$의 이론적 기여를 실무적 가치로 전환하기 위한 구체적 기반을 제공하며,
\tmod{\chref{ch:discussion}}{\tsecref{sec:conclusion-future-merged}}에서는 이러한 논의를 종합하여 본 연구의 결론과 향후 연구 방향을 제시한다.% 3차 심사 (2026-05-09): 자기참조 오류 수정 — \chref{ch:discussion}가 이 5장 자체를 가리킴. 5.5절을 가리키도록 \tsecref(빨간색 secref)로 수정.


% --- 5.1.3에서 흡수: 종합 논의 절에 있던 ``연구의 한계'' 4단락 ---

%% -------------------------------------------------------------------
\subsection{\syncr{연구의 한계 (소규모 패널 \& 단일 $\kappa$)}}
\label{subsec:ch5_limitations}
%% -------------------------------------------------------------------

본 연구의 한계는 다음과 같다.

\fdel{\noindent\textbf{시뮬레이션 기반 검증의 한계.}
본 연구의 결과는 $n = 10{,}000$ 몬테카를로 시뮬레이션에 기반하며,
실제 장애 사건 데이터에 의한 검증은 수행되지 않았다.
시뮬레이션에서 사용한 입력 분포($k_O, O \sim \mathrm{Uniform}(0,1)$,
$U \sim N(3,1)$ 또는 $6 \cdot \mathrm{Pareto}(3)$)는 이론적
분석에 적합하나, 실제 조직의 관측성과 불확실성 분포를
완전히 반영하지 못할 수 있다.
향후 산업별 실제 장애 이력 데이터를 활용한 경험적 검증이 필요하다.}

\chg{전문가 패널의 규모 측면에서,}
\chg{정량적과 정성적} 전문가 검증을 5명의 소규모 패널로 수행하였다.
S-CVI/Ave $= 0.985$, Cronbach's $\alpha = 0.89$로 내용 타당도 및 내적 일관성이
확인되었으나, 대규모 패널에 의한 일반화 가능성은 추가 검증이 필요하다.
특히 산업별, 규모별 조직의 다양성을 반영한 전문가 패널 확장이
후속 연구에서 수행되어야 한다.

\chg{단일 $\kappa$ 값의 적용 측면에서,}
본 연구는 \jrev{\chg{\citet{lee2026drto}}에서 \chg{도출한}} $\kappa = 1.2$의 단일 곡률 매개변수를 전체 조건에 적용하였다.
\jrev{본 연구는 \frev{분포 환경 적합성} 종합 명제(\secref{sec:ch4_hypothesis_summary})를 통해 9개 가설(H2--H10) PASS의 통합 해석으로 $\kappa = 1.2$의 C2/C3 조건 \frev{운영 적합성을 시사하였으나},}
$\kappa$의 조건별 또는 산업별 차별화\jrev{, 그리고 환경 무관 다기준 종합점수의 재정의를 통한 직접 격자 탐색}이 모델의 적응력을 향상시킬 수 있으며,
이는 후속 연구의 과제이다.

\fdel{\noindent\textbf{시간적 역동성의 미반영.}
현재 $\DRTO$ 모델은 특정 시점의 관측성과 불확실성을 입력으로 하는
정적(static) 산출 구조이다. 실제 복구 과정에서 관측성과 불확실성은
시간에 따라 변화하므로, 시간적 역동성을 반영하는
동적(dynamic) 갱신 모델로의 확장이 필요하다.}

    % 5.4 한계점 (+5.1.3 흡수)
% =====================================================================
% 5.5 결론 및 향후 연구 — 5단계 5-2 (clean subsection wrapper)
% 본문 텍스트 무수정. 분할 본문 파일을 3항 구조로 재배치.
%   5.5.1 연구 요약       ← 현 sec5.5.1_summary.tex
%   5.5.2 주요 연구 성과  ← 현 sec5.5.2_achievements.tex (5.7.x → 5.5.2.x로 강등)
%   5.5.3 향후 연구 방향  ← 현 sec5.5.3_future_work.tex
% =====================================================================

\section{\vrev{결론 및 향후 연구}}
\label{sec:conclusion-future-merged}

\begin{p5add}
본 절에서는 본 연구의 전체 흐름을 압축
요약하고(5.5.1항), 핵심 성과 5개를 정리하며(5.5.2항), 한계점에서 도출된
향후 연구 방향(5.5.3항)을 제시한다. 이로써 본 연구의 학술적 기여
범위와 실무적 적용 가능성을 명확히 하고, 후속 연구의 출발점을 마련한다.
\end{p5add}

% --- 5.5.1 연구 요약 ---
\subsection{\vrev{연구 요약}}
\label{subsec:conclusion-summary-merged}

\begingroup
  \renewcommand{\section}[1]{}
  %% ===================================================================
\section{\vrev{연구 요약}}
\label{sec:conclusion-summary}
%% ===================================================================

현행 BCMS에서 RTO는 BIA 시점에
고정값으로 설정되며, 이후 운영 환경의 변화와 무관하게 유지된다.
이러한 정적(static) RTO 접근법은 세 가지 구조적 한계를 내포한다.
첫째, 조직의 관측성(observability) 수준이 시간과 상황에 따라 변동함에도
이를 반영하지 못한다.
둘째, 복구 과정에 내재된 불확실성(uncertainty)을 확률 분포에 기반하여
정량화하지 않는다.
셋째, 사이버 공격이나 복합 재난 등 극단적 사건의 파레토(중꼬리, Heavy-Tail) 특성을
구조적으로 포착하지 못한다.

본 연구는 이러한 한계를 극복하기 위하여, 관측성과 불확실성을
개별 시그모이드 바운딩(individual sigmoid bounding) 함수로 통합하는
$\DRTO$ 프레임워크를 제안하였다.
핵심 수식은 \m{식~(\ref{eq:ch7-drto-core})에 나타낸 바와 같이} 정의된다.

\begin{equation}
  \DRTO = R_0 + E_{O,\text{bnd}} + E_{U,\text{bnd}}
  \label{eq:ch7-drto-core}
\end{equation}

\noindent
여기서 관측성 효과 $E_{O,\text{bnd}} = -R_0 \cdot S(z_O)$는
$\DRTO$를 기준선 $R_0$로부터 하향 조정하고,
불확실성 효과 $E_{U,\text{bnd}} = (R_{\max} - R_0) \cdot S(z_U)$는
상향 조정하며, 수정 시그모이드 $S(z) = 2/(1 + e^{-z}) - 1$에 의해
물리적 경계 $\DRTO \in (0,\; R_{\max})$가 점근적으로 보장된다.
$R_0 = 6.0$시간, $R_{\max} = 24.0$시간, $\kappa = 1.2$를 기본 설정으로 채택하였다.

연구 검증은 세 축으로 수행되었다.
첫째, $n = 10{,}000$ 몬테카를로 시뮬레이션을 4개 조건(C0--C3)에 대해
실시하여 \jrev{10개 연구가설}(H1--\jrev{H10})의 모델 내부 정합성을 체계적으로 확인하였다.
둘째, BCM/DR 분야 전문가 5명을 대상으로 구조화된 설문(27문항, 5점 리커트 척도)과
심층 인터뷰(약 30분, 반구조화)를 수행하여 모델의 실무 타당성을 삼각 검증하였다.
셋째, 1차 단순, 1차 다중, 2차 회귀 분석 및 P85.8 기준 분할 회귀를 통해
모델 응답 곡면의 구조적 특성을 규명하고, 꼬리 감쇠 현상을 발견하였다.
그 결과, \textbf{\jrev{10개} 가설 전원이 사전 설정 기준을 충족하여 채택}되었다.
넷째, \chg{\citet{lee2026drto}}과 본 연구의 5개 연구 단계에 대한
\chg{연구모형, 온톨로지 분석, 논문 내용, 기준값} 간
\textbf{303개 교차검증 항목을 전수 확인}하여
\chg{수치, 수식, 용어}의 완전 정합(전원 PASS)을 달성하였으며,
이는 연구 프레임워크의 논리적 일관성과 데이터 신뢰성을 보증한다.



\endgroup

% --- 5.5.2 주요 연구 성과 ---
\subsection{\vrev{주요 연구 성과}}
\label{subsec:conclusion-achievements-merged}

\begingroup
  \renewcommand{\section}[1]{}
  \renewcommand{\subsection}[1]{\subsubsection{#1}}
  %% ===================================================================
\section{\vrev{주요 연구 성과}}
\label{sec:conclusion-achievements}
%% ===================================================================

\revdel{본 연구의 5대 핵심 성과를 다음과 같이 정리한다.}\,\revadd{본 항은 5장 전체에서 도출된 5대 핵심 성과를 압축 요약 형태로 제시한다. 각 성과의 상세 논의는 학술적 측면은 \secref{sec:academic-contrib}, 실무적 측면은 \secref{sec:practical-contrib}을 참조하라.}


%% -------------------------------------------------------------------
\chg{첫째 성과는 물리적 경계 보장을 갖춘 개별 시그모이드 바운딩 $\DRTO$ 모델이다.}
\label{subsec:achiev-model}
\chg{개별 시그모이드 바운딩을 통해 $\DRTO \in (0, R_{\max})$를 보장하고, $C^{\infty}$ 연속성과 자연스러운 교호작용 포착을 함께 달성하였다(상세는 \subsecref{subsec:academic-theoretical}).}

\fdel{업무연속성관리 분야 최초로, 관측성과 불확실성을 단일 수리 모델에서
통합하는 개별 시그모이드 바운딩 $\DRTO$ 모델을 정립하였다.
수정 시그모이드 함수 $S(z) = 2/(1 + e^{-z}) - 1$을 기저 함수로 채택하여,
관측성 효과($E_{O,\text{bnd}}$)와 불확실성 효과($E_{U,\text{bnd}}$)를
각각 고유한 물리적 범위($(-R_0,\; 0]$과 $[0,\; R_{\max} - R_0)$)에서
독립적으로 바운딩하였다. 이 모델은 물리적 경계 보장($\DRTO \in (0, R_{\max})$),
무한 미분 가능성($C^{\infty}$), 교호작용의 자연스러운 포착이라는
세 요구사항을 동시에 충족한다.
기존의 클리핑(clipping) 기반 접근이나 경계 미보장 모델과 달리,
시그모이드 바운딩은 점근적으로 경계에 도달하되 경계값 자체를 취하지 않아
수치적 안정성과 이론적 엄밀성을 동시에 확보한다.}


%% -------------------------------------------------------------------
\chg{둘째 성과는 \citet{lee2026drto}의 $\kappa^{*} = 1.2$에 대한 \frev{분포 환경 적합성} 종합 명제 도출이다.}
\label{subsec:achiev-kappa}
\chg{\citet{lee2026drto}의 $\kappa^{*}=1.2$를 본 연구의 새로운 분포 환경(C2 가우시안과 C3 파레토)에 적용한 결과, 9개 가설(H2--H10) PASS의 통합 해석을 통해 $\kappa = 1.2$의 분포 환경 \frev{운영 적합성}을 시사하였다(상세는 \secref{sec:ch4_hypothesis_summary}의 \frev{분포 환경 적합성} 종합 명제와 \subsecref{subsec:academic-methodological}).}

\fdel{시그모이드 곡률 파라미터 $\kappa$의 최적값을 다중 기준 복합 점수(multi-criteria
composite score) 방법론으로 도출하였다.
민감도(sensitivity), 선형성(linearity), 경계 여유(boundary margin),
해석 가능성(interpretability)의 4개 기준을 정규화하여 종합 평가한 결과,
$\kappa^{*} = 1.2$가 최적점으로 결정되었다.
$\kappa = 1.2$에서 관측성 채널의 $z_O \in [0,\; 1.2]$는
시그모이드의 비선형 압축이 의미 있게 작동하는 영역에 해당하여
($S(1.2) = 0.537$), 중간 입력에서의 감도를 확보하면서도
극단 영역에서의 비선형 포화(saturation) 거동이 물리적 경계 보장을
자연스럽게 실현한다.}


%% -------------------------------------------------------------------
\chg{셋째 성과는 C0$\to$C3 4단계 점진적 프레임워크와 \jrev{10개} 가설의 전원 채택이다.}
\label{subsec:achiev-framework}
\chg{C0--C3 4단계 점진적 조건과 6단계 다층 검증을 통해 \jrev{10개} 가설(H1--\jrev{H10})을 전원 채택하였다(상세는 \subsecref{subsec:academic-empirical}).}

\fdel{정적 기준선(C0), 관측성 도입(C1), 가우시안 불확실성(C2),
파레토 불확실성(C3)의 4단계 점진적 조건 체계를 수립하고,
6단계 다층 검증 프레임워크(L0 기반 검증(Golden Compare),
시뮬레이션 검증(H1--\jrev{H3}), 회귀 검증(\jrev{H4}--\jrev{H5}), 강건성 검증(\jrev{H6}--\jrev{H7}),
이중채널 검증(\jrev{H8}), 전문가 평가(\jrev{H9}--\jrev{H10}))를 통해 \jrev{10개} 가설을 체계적으로 검증하였다.
정량적 결과의 핵심은 다음과 같다.
C1에서 관측성에 의한 $\eta$ 감소($\bar{\eta} = 0.853$, Cohen's $d = -1.180$),
C2에서 불확실성 프리미엄의 발현($\bar{\eta} = 1.682$, Cohen's $d = +1.871$),
C3에서 평균 유사$\cdot$꼬리 급증($\bar{\eta} = 1.514$, CVaR$_{99}$ $+24.0\%$).
$R^2_{(2\mathrm{q})}$는 C1$\,= 1.000$, C2$\,= 0.998$이며,
C3는 전체 $R^2_{(2\mathrm{q})} = 0.858$이나 P85.8 분할 후 Core $R^2 = 0.999$로 \jrev{H4}가 지지되었다.}


%% -------------------------------------------------------------------
\chg{넷째 성과는 꼬리 감쇠 현상의 발견($0.52$배, 꼬리 영역)이다.}
\label{subsec:achiev-dual-channel}
\chg{P85.8 분할 회귀에서 $k_O$ 계수가 Core $-0.798$에서 Tail $-0.415$로 $0.52$배 감쇠하였고, CVaR$_{99}$는 $+24.0\%$ 증가하였다(상세는 \subsecref{subsec:academic-theoretical}).}

\fdel{C3 조건에서 P85.8 기준으로 분할 회귀를 수행한 결과,
관측성 가중치 $k_O$의 회귀 계수가 핵심 영역(Core, $-0.798$)에서
꼬리 영역(Tail, $-0.415$)으로 $0.52$배 감쇠되는 꼬리 감쇠
(tail attenuation) 현상을 발견하였다.
이 발견은 극단적 불확실성 환경에서 시그모이드 포화에 의해
관측성 투자의 한계 효용이 완화된다는 이론적 시사점을 제공한다.
P85.8 CDF 교차점은 가우시안(C2)과 파레토(C3)의 CDF가
교차하는 구조적 경계로서, 교차점 이하에서 두 분포는
사실상 구분 불가능하나(Cohen's $d \approx 0$),
교차점 초과에서 C3의 $\eta$가 C2를 급격히 상회하여
CVaR$_{99}$ 기준 $+24.0\%$의 격차를 보인다.}


%% -------------------------------------------------------------------
\chg{다섯째 성과는 전문가 검증(4.46/5.0)과 6단계 다층 검증이다.}
\label{subsec:achiev-validation}
\chg{전문가 5명의 평균이 4.46/5.0($p=0.011$, $d=1.65$)이고 Cronbach's $\alpha=0.89$, S-CVI/Ave $=0.985$였으며, 인터뷰에서 15개 주제 모두 $\geq 80\%$ 동의를 얻었다(상세는 \subsecref{subsec:academic-empirical}).}

\fdel{BCM/DR 분야 전문가 5명(평균 경력 15.2년)을 대상으로
수행한 정량적 설문에서 전체 평균 4.46/5.0점($t(4) = 3.68$, $p = 0.011$,
Cohen's $d = 1.65$)을 달성하였으며,
Cronbach's $\alpha = 0.89$, S-CVI/Ave$\,= 0.985$로
높은 내적 일관성과 평가자 간 일치도가 확인되었다(\jrev{H9} 지지).
심층 인터뷰에서는 5대 주제 15개 하위 주제 모두(100\%)에서 80\% 이상
전문가 동의가 달성되어 \jrev{H10}이 지지되었다.
시뮬레이션 검증, 회귀 검증, 강건성 검증, 이중채널 검증, 전문가 평가의
6단계(L0--L5)가 계층적 의존 구조를 형성하여 모델의 타당성을 다각적으로 보장하며,
정량적 시뮬레이션과 정성적 전문가 평가의 삼각검증(triangulation)을 통해
방법론적 편향을 최소화하였다.}

\vspace{0.5em}

\p{[}표~\vref{tab:ch7-findings-summary}\p{]}는 이상의 5대 성과를 핵심 결과,
정량적 근거, 학술적 의의의 관점에서 종합한 것이다.

\begin{table}[htbp]
\centering
\caption{핵심 연구 결과 요약}
\label{tab:ch7-findings-summary}
\small
\begin{tabularx}{\textwidth}{l X X}
\toprule
\textbf{핵심 결과} & \textbf{정량적 근거} & \textbf{학술적 의의} \\
\midrule
개별 시그모이드 바운딩 모델
  & $\DRTO \in (0,\; R_{\max})$ 보장, $C^{\infty}$ 연속
  & BCM 분야 최초의 물리적 경계 보장 동적 RTO 모델 \\
\addlinespace
\jrev{$\kappa^{*}=1.2$ \frev{분포 환경 적합성}}
  & \jrev{9개 가설(H2--H10) PASS 통합 해석}
  & \jrev{\chg{\citet{lee2026drto}의} 매개변수의 본 연구 환경 적용성을 다중 가설 PASS로 \frev{시사한} 학술적 명확성} \\
\addlinespace
C0$\to$C3 프레임워크, H1--\jrev{H10} 전원 채택
  & $\eta$: 0.853(C1), 1.682(C2), 1.514(C3); C1/C2 $R^2 \geq 0.998$, C3 Core $R^2 = 0.999$
  & 4단계 점진적 BCM 성숙도 모형의 실증적 근거 \\
\addlinespace
꼬리 감쇠 ($0.52$배)
  & Tail $k_O$ 계수 $-0.415$ vs Core $-0.798$; CVaR$_{99}$ $+24.0\%$
  & 극단 환경에서 관측성 투자 효용 감소의 최초 정량적 규명 \\
\addlinespace
전문가 검증 + 6단계 다층 검증
  & 4.46/5.0, $\alpha = 0.89$, S-CVI/Ave$\,= 0.985$; \jrev{10/10} 가설 채택
  & \chg{시뮬레이션과 전문가} 삼각검증의 방법론적 표준 제시 \\
\bottomrule
\end{tabularx}
\end{table}

\chg{이상의 5대 성과는 학술적 의의에 그치지 않고 본 장의 실무적 기여로 직접 이어진다. 개별 시그모이드
바운딩 모델과 회귀 설명력(첫째와 셋째 성과)은 회귀식 기반 $\DRTO$ 산출 도구(\subsecref{subsec:practical-regression-tool})로,
꼬리 감쇠 발견(넷째 성과)은 이중채널 감쇠 기반 자원배분과 스트레스 테스트
도구(\subsecref{subsec:practical-dualchannel-tool})로, $\kappa = 1.2$의 적합성과 전문가 검증(둘째와 다섯째 성과)은
산업별 적용 가이드라인과 운영 도구의 신뢰 근거로 활용된다. 한편 위 표가 다루는 다섯 핵심 성과 외에도
4장에서 도출된 세부 결과들이 의미의 크고 작음과 무관하게 각각 어느 학술 기여 또는 실무 활용으로
귀착되는지는 \subsecref{subsec:ch5_results_map}의 전수 매핑에 정리되어 있다.}



\endgroup

% --- 5.5.3 향후 연구 방향 ---
\subsection{\vrev{향후 연구 방향}}
\label{subsec:conclusion-future-merged-future}

\begingroup
  \renewcommand{\section}[1]{}
  %% ===================================================================
\section{\vrev{향후 연구 방향}}
\label{sec:conclusion-future}
%% ===================================================================

본 연구의 한계를 토대로 \jrev{5개} 향후 연구 방향을 다음과 같이 제안한다.


\chg{첫째, 실제 조직 데이터 기반 현장 검증이다.}
본 연구는 시뮬레이션 기반 검증에 초점을 맞추었으며, 실제 조직 환경에서의
현장 실험(field test)을 수반하지 않았다.
제조업, 금융업, IT 서비스 등 산업별 대표 조직을 선정하여
실제 장애 이벤트에서의 복구 시간을 수집하고,
시뮬레이션에서 예측한 $\eta$ 값과의 일치도를 검증함으로써
모델의 외적 타당성(external validity)을 확보해야 한다.
특히 관측성 지표($O(t)$)의 실시간 수집 방법론과
불확실성 추정을 위한 장애 이력 데이터베이스 구축 방안을 구체화하여,
$\DRTO$의 실무적 운용 기반을 마련할 필요가 있다.


\chg{둘째, 시간 동역학(temporal dynamics) 및 상태공간 모형으로의 확장이다.}
현행 $\DRTO$ 모델에서 $R_0$, $R_{\max}$, $\kappa$는 고정값으로 설정되어 있다.
이를 시간 변동 함수\chg{(}$R_0(t)$, $\kappa(t)$\chg{)}로 확장하여
조직의 BCM 성숙도 변화와 환경 변동을 모델에 내생적으로 반영할 수 있다.
상태 공간 모형(state-space model)의 관점에서 $\DRTO$의 시간적 진화를
형식화하면, 관측성과 불확실성의 동적 상호 작용을
시간축 위에서 추적할 수 있는 이론적 기반이 마련된다.
특히 이산 시간 상태 전이 행렬을 통해
C0$\to$C1$\to$C2(가우시안) 및 C0$\to$C1$\to$C3(파레토)의
분기 경로를 동적 체제 전환(dynamic regime switching)~\citep{hamilton1989regime}으로
모형화하는 확장이 가능하다.


\chg{셋째, 칼만 필터(Kalman filter) 통합이다.}
상태 공간 모형의 자연스러운 확장으로서, 칼만 필터를 $\DRTO$ 프레임워크에
통합하여 입력 변수($O(t)$, $U(t)$)의 노이즈를 실시간으로 여과하고
최적 상태 추정을 수행할 수 있다.
장애 이벤트 관측 시마다 예측-갱신(\erf{pict-update}{predict-update}) 사이클을 통해
$\DRTO$를 적응적으로 교정하는 자기 교정(self-calibrating) 메커니즘은,
관측성 지표의 측정 오차와 불확실성 추정의 편향을 체계적으로 보정할 수 있다.
확장 칼만 필터(EKF) 또는 언센티드 칼만 필터(UKF)를 적용하면
시그모이드 바운딩의 비선형성을 직접 처리할 수 있으며,
이는 $\DRTO$의 실시간 운용 정확도를 향상시킬 것이다.


\chg{넷째, \frev{분포 환경 적합성}의 \frev{직접 검증}이다.}
\jrev{본 연구는 \chg{\citet{lee2026drto}}의 $\kappa^{*}=1.2$를 새로운 분포 환경에 적용한 후 9개 가설(H2--H10) PASS의 통합 해석으로 \frev{운영 적합성을 시사하였다} (4.13절 \frev{분포 환경 적합성} 종합 명제). 그러나 \chg{\citet{lee2026drto}}의 4기준 종합 점수 $F(\kappa) = f_1 \cdot f_2 \cdot f_3 \cdot f_4$ (\chg{실무 유의성, 비포화, 효과 크기, 설명력})는 C1 조건 기반으로 정의되어 C2/C3 조건에서 직접 격자 탐색에 부적합하다는 학술적 한계를 내포한다. 향후 연구에서는 분포 환경 무관 다기준 종합 점수의 재정의와 C2/C3 조건에서의 직접 격자 탐색을 통해 $\kappa^{*}$의 환경별 최적값 비교를 수행함으로써, 본 연구의 종합 명제를 직접적 검증으로 보완할 수 있다.}


\chg{다섯째, 다중 조직 BCM 및 적응적 $\kappa$ 강화학습이다.}
단일 조직의 $\DRTO$에서 확장하여, 복수의 조직이 상호 의존적 공급망을
형성하는 다중 조직 수준의 $\DRTO$ 모형을 개발할 수 있다.
시스템 간 장애 전파(cascading failure) 효과, 공유 자원 경쟁,
복구 우선순위 최적화를 개별 시그모이드 바운딩의 다차원 확장을 통해 모형화한다.
나아가 강화 학습(reinforcement learning)을 적용하여
$\kappa$를 환경 변화에 따라 적응적으로 조정하는 메커니즘을 구축할 수 있다.
조직의 장애 이력과 복구 성과를 보상(reward) 신호로 활용하면,
$\kappa$가 \chg{조직과 산업, 그리고 시점}에 최적화된 값으로
자동 수렴하는 자기 최적화(self-optimizing) 프레임워크의 실현이 가능하다.

\vspace{0.5em}

\p{[}그림~\vref{fig:ch7-achievement-map}\p{]}는 본 연구의 연구 성과 구조를
문제 인식에서 해결 방안, 핵심 성과로 이어지는 흐름으로 시각화한 것이다.

\begin{figure}[htbp]
\centering
\resizebox{\textwidth}{!}{%
\begin{tikzpicture}[
    probbox/.style={draw, rounded corners=5pt, fill=red!8, minimum width=34mm,
                    minimum height=12mm, align=center, font=\footnotesize,
                    line width=0.6pt},
    solbox/.style={draw, rounded corners=5pt, fill=blue!8, minimum width=34mm,
                   minimum height=12mm, align=center, font=\footnotesize,
                   line width=0.6pt},
    outbox/.style={draw, rounded corners=5pt, fill=green!8, minimum width=34mm,
                   minimum height=12mm, align=center, font=\footnotesize,
                   line width=0.6pt},
    catbox/.style={draw, rounded corners=8pt, inner sep=6pt, line width=0.8pt},
    conn/.style={-{Stealth[length=2.5mm]}, thick, black},
    catlbl/.style={font=\small\bfseries, anchor=south},
]

% --- 문제 인식 (왼쪽) ---
\node[probbox] (p1) at (0, 3.0) {관측성 미반영\\(정적 RTO)};
\node[probbox] (p2) at (0, 1.5) {불확실성 미정량화\\(worst-case 한정)};
\node[probbox] (p3) at (0, 0.0) {극단 사건 미모형화\\(정규분포 전제)};

% --- 해결 방안 (중앙) ---
\node[solbox] (s1) at (5.5, 3.0) {개별 시그모이드\\바운딩 모델};
\node[solbox] (s2) at (5.5, 1.5) {\frev{분포 환경 적합성}\\종합 명제 도출};
\node[solbox] (s3) at (5.5, 0.0) {C0$\to$C3 점진적\\4단계 프레임워크};

% --- 핵심 성과 (오른쪽) ---
\node[outbox] (o1) at (11, 3.6) {$\DRTO \in (0, R_{\max})$\\경계 보장};
\node[outbox] (o2) at (11, 2.1) {H1--H10 전원 채택\\$R^2_{(2\mathrm{q})} \geq 0.95$};
\node[outbox] (o3) at (11, 0.6) {꼬리 감쇠 $0.52$배\\CVaR$_{99}$ $+24.0\%$};
\node[outbox] (o4) at (11, -0.9) {전문가 4.46/5.0\\$\alpha = 0.89$};

% --- 카테고리 박스 (배경 레이어) ---
\begin{scope}[on background layer]
  \node[catbox, fit=(p1)(p2)(p3), draw=red!50, fill=red!2] (pcat) {};
  \node[catbox, fit=(s1)(s2)(s3), draw=blue!50, fill=blue!2] (scat) {};
  \node[catbox, fit=(o1)(o2)(o3)(o4), draw=green!50!black, fill=green!2] (ocat) {};
\end{scope}

% --- 카테고리 라벨 ---
\node[catlbl, text=black] at (pcat.north) {문제 인식};
\node[catlbl, text=black] at (scat.north) {해결 방안};
\node[catlbl, text=black] at (ocat.north) {핵심 성과};

% --- 연결선 ---
\draw[conn] (p1.east) -- (s1.west);
\draw[conn] (p2.east) -- (s2.west);
\draw[conn] (p3.east) -- (s3.west);

\draw[conn] (s1.east) -- (o1.west);
\draw[conn] (s2.east) -- (o2.west);
\draw[conn] (s3.east) -- (o3.west);
\draw[conn, dashed] (s1.east) -- (o2.west);
\draw[conn, dashed] (s3.east) -- (o4.west);

\end{tikzpicture}%
}% end resizebox
\caption{문제 인식에서 해결 방안, 핵심 성과로 이어지는 연구 성과 구조도}
\label{fig:ch7-achievement-map}
\end{figure}


\vspace{1.5em}

결론적으로, 본 연구는 \erf{BCMgownj}{BCM} 분야에서 정적 RTO의 근본적 한계를
극복하는 $\DRTO$ 프레임워크를 제안하고,
개별 시그모이드 바운딩 모델\jrev{(\chg{\citet{lee2026drto}}의 $\kappa^{*} = 1.2$ 전제 채택)},
꼬리 감쇠 발견($0.52$배), P85.8 CDF 교차점 규명(CVaR$_{99}$ $+24.0\%$),
6단계 다층 검증(전문가 4.46/5.0, \jrev{10/10} 가설 채택)\jrev{,
그리고 \textbf{\jrev{\frev{분포 환경 적합성}} 종합 명제 도출}}이라는
\chg{이론적이고 방법론적인} 기여를 달성하였다.
향후 실제 조직 데이터 기반 현장 검증, 상태 공간 모형과 칼만 필터의 통합,
\jrev{\frev{분포 환경 적합성}의 \frev{직접 검증}(환경 무관 다기준 종합 점수 재정의),}
다중 조직 BCM 확장, 그리고 적응적 $\kappa$ 강화 학습을 통해
$\DRTO$ 프레임워크가 조직의 업무연속성 역량을 실질적으로 강화하는
정량적 도구로 발전하기를 기대한다.


\endgroup
             % 5.5 결론 및 향후 연구
